\documentclass[10pt,a4paper]{article}

% ----------------- Paquetes -------------------%
\usepackage[T1]{fontenc}
\usepackage[utf8]{inputenc}
\usepackage[a4paper,margin=2.5cm]{geometry}
\usepackage{mathptmx}             
\usepackage{changepage} 
\usepackage{setspace}
\usepackage[spanish]{babel}
\usepackage[labelfont=bf,labelsep=space]{caption}
\usepackage{amssymb}
\usepackage{amsmath}
\usepackage{ebgaramond}
\usepackage{placeins}
\usepackage{etoolbox}
\usepackage{setspace}
\usepackage{parskip} 
\usepackage{tcolorbox}
\usepackage{needspace}
\usepackage{xparse}
\usepackage{ocgx2}
\usepackage{hyperref}
\usepackage{hypcap}
\usepackage{soul}\sethlcolor{yellow}% resaltado amarillo: \hl{texto} (Panel TCEE, ⌃H)


%%%%%%%%%%%%%%%%%%%%%%%%%%%%%%%%%%%%%%%%%%%%%%%%%%%%%%%%%%%%%%%%
% --- Fija primero tus valores globales "buenos" ---
\setstretch{1.2}                 % Interlineado global
\setlength{\parskip}{0.7\baselineskip}
\setlength{\parindent}{0pt}
\newlength{\GlobalParskip}
\setlength{\GlobalParskip}{\parskip}

\newcounter{eqnum}[subsection]
\renewcommand{\theeqnum}{\arabic{eqnum}}
\newcounter{alphaitem}
\newcounter{numitem}

%% ------------------- Recuadros----------------------%%
\newcommand{\puntosclave}[1]{%
  \begin{tcolorbox}[
    colframe=black,
    title={Puntos clave},
    before upper={\setlength{\parindent}{0.5cm}\setlength{\parskip}{0.5\baselineskip}}
  ]
  #1
  \end{tcolorbox}
}

\newcommand{\modificaciones}[1]{
  \begin{tcolorbox}[
    colback=white,
    colframe=red,
    title={Modificaciones a realizar},
    before upper={\setlength{\parindent}{0.5cm}\setlength{\parskip}{0.5\baselineskip}}
  ]
  #1
  \end{tcolorbox}
}

%% ------------------- CITA ----------------------%%
% \begin{cita}[Autor][Año][Obra] texto \end{cita}: cita sangrada y, debajo a la derecha, «AUTOR (Año), Obra». Los tres datos son opcionales.
% En el Panel TCEE: escribir \cita e Intro (el cursor va al texto; Tab pasa a Autor, Año y Obra).
\NewDocumentEnvironment{cita}{ O{} O{} O{} +b }{%
  \begin{adjustwidth}{2cm}{0pt}
  \raggedright
  #4\par
  \raggedleft
  \if\relax\detokenize{#1#2#3}\relax
    % sin firma
  \else
    #1%
    \if\relax\detokenize{#2}\relax\else\if\relax\detokenize{#1}\relax\else\space\fi(#2)\fi
    \if\relax\detokenize{#3}\relax\else\if\relax\detokenize{#1#2}\relax\else, \fi\textit{#3}\fi
  \fi
  \end{adjustwidth}
}{}



%% ------------------- Listas----------------------%%
    % --- Lista alfabética ---
    \newenvironment{la}
      {\begin{list}{\alph{alphaitem})}{%
          \usecounter{alphaitem}%
          \setlength{\leftmargin}{1cm}%
          \setlength{\parskip}{3pt}% 
          \setlength{\itemsep}{0pt}% ← espacio entre ítems
          \setlength{\parsep}{0pt}%  ← párrafos dentro de un ítem
      }}
      {\end{list}}
      
    % --- Lista numérica ---
    \newenvironment{lnum}
      {\begin{list}{\arabic{numitem}.}{%
          \usecounter{numitem}%
          \setlength{\leftmargin}{1cm}%
          \setlength{\parskip}{3pt}% 
          \setlength{\itemsep}{0pt}% ← espacio entre ítems
          \setlength{\parsep}{0pt}%  ← párrafos dentro de un ítem
      }}
      {\end{list}}
      
% ------------------- Imágenes----------------------%
    \graphicspath{{Img/}}% Rutas por defecto 
    % --- Imagen adaptativa ---
    % Registros de dimensión definidos una sola vez
    \newdimen\imgcap
    \newdimen\imgmin
    \newdimen\imgmax
    \newdimen\imgremain
    \newdimen\imgh
    \newdimen\imgneed
    
    \newcommand{\imagenfit}[3]{%
      \addvspace{10pt}% separación antes
      \par\begingroup
        % Parámetros de composición (asignaciones, no \newdimen)
        \imgcap=1\baselineskip            % alto estimado de título+fuente
        \imgmin=0.15\textheight
        \imgmax=0.15\textheight
        \imgremain=\pagegoal
        \ifdim\pagegoal=\maxdimen \imgremain=0pt\fi
        \advance\imgremain by -\pagetotal
        \advance\imgremain by -\imgcap
        % Decisión de altura destino
        \ifdim\imgremain<\imgmin
          \newpage
          \imgh=0.20\textheight
        \else
          \imgh=\imgremain
          \ifdim\imgh>\imgmax \imgh=\imgmax \fi
        \fi
        % Reservar espacio para evitar cortes del bloque completo
        \imgneed=\imgh \advance\imgneed by \imgcap
        \Needspace{\imgneed}
        % Cuerpo
        \noindent\begin{minipage}{\textwidth}
          \centering
          \captionof{figure}{\textemdash\ #2}\par\vspace{0.3em}% título arriba
          \includegraphics[width=\linewidth,height=\imgh,keepaspectratio]{\detokenize{#1}}\par
          \vspace{0.3em}\small\textit{Fuente: }#3% fuente abajo
        \end{minipage}%
      \endgroup
      \par\addvspace{10pt}%
    }

%------------ Bloque de RECUADRO ESPECIAL -------------%
    \newenvironment{specialblock}[1]{%
      \begin{quote} % crea sangría a izquierda y derecha
      \small % texto más pequeño
      \noindent\textbf{#1}\\[-0.3em] % título en negrita
      \rule{\linewidth}{0.4pt}\\[0.6em]}{
      \end{quote}}
      

%-------------------------- Author Font -----------------------%
    \newcommand{\authorfont}[1]{{\fontfamily{EBGaramond-TLF}\selectfont #1}}

%-------------------------- Anexos -----------------------%
    \makeatletter
    \newcounter{anexo}
    \renewcommand{\theanexo}{\Roman{anexo}}
    \newcommand{\anexo}[1]{%
      \clearpage
      \phantomsection
      \refstepcounter{anexo}%
      \noindent\textbf{Anexo \theanexo.- }#1\par\medskip
      \addcontentsline{stc}{section}{Anexo \theanexo.- #1}%
    }
    \makeatother

    %-------------------------- Lista de figuras (personal) -----------------------%
\makeatletter
\newcommand{\MyListOfFigures}{%
  \clearpage
  \phantomsection
  {\centering\bfseries\Large Índice de figuras\par}%
  \medskip
  % Entrada en nuestro índice de contenido (stc)
  \addcontentsline{stc}{section}{Índice de figuras}%
  % Recuperación del listado (sin cabecera propia)
  \begingroup
    \parindent=0pt\parskip=2pt
    \@starttoc{lof}%
  \endgroup
}
\makeatother

%-------------------------- Bibliografía -----------------------%
\makeatletter
% Contador para las referencias
\newcounter{myref}
% Cabecera de la bibliografía, entra en el índice 'stc'
\newcommand{\MyBibliography}{%
  \clearpage
  \phantomsection
  {\centering\bfseries\Large Bibliografía y referencias\par}%
  \medskip
  \addcontentsline{stc}{section}{Bibliografía y referencias}%
  \setcounter{myref}{0}%
}
\newcommand{\MyBibItem}[4]{%
  \refstepcounter{myref}%
  \noindent[\themyref]\ #1.\ \emph{#2}. #3. #4\par
}
\makeatother

%--------- Establecer cuatro niveles de índice --------%
    \setcounter{tocdepth}{4}   
    \setcounter{secnumdepth}{4}% numera hasta \paragraph

%%%%%%%%%%%%%%%%%%%%%%%%%%%%%%%%

\makeatletter

    %--------- Interlineado Global --------%
    \edef\GlobalBaseLineStretch{\baselinestretch}
    
    %--------- Espaciado de seccinones --------%
    \renewcommand\section{\@startsection{section}
      {1}{\z@}
      {2.5ex \@plus 1ex \@minus .2ex} % espacio antes
      {1.5ex \@plus .2ex}             % espacio después
      {\normalfont\Large\bfseries}    % formato del título
    }
    \renewcommand\subsection{\@startsection{subsection}{2}{\z@}
      {3ex \@plus 0.8ex \@minus .2ex}
      {1.5ex \@plus .2ex}
      {\normalfont\large\bfseries}
    }
    \renewcommand\subsubsection{\@startsection{subsubsection}{3}{\z@}
      {3ex \@plus 0.5ex \@minus .2ex}
      {1ex \@plus .2ex}
      {\normalfont\normalsize\bfseries}
    }
    \renewcommand\paragraph{\@startsection{paragraph}{4}{\z@}%
    {-2ex \@plus -0.5ex \@minus -.2ex}% espacio antes
    {0.8ex \@plus .2ex}% espacio después
    {\normalfont\normalsize\itshape}}% ← cursiva en lugar de negrita

    %--------- Ecuaciones --------%   
    % Variables globales para almacenar títulos y notas
    \gdef\leq@current@section{}%
    \gdef\leq@current@subsection{}%
    \gdef\leq@last@printed@section{}%
    \gdef\leq@last@printed@subsection{}%
    
    % Comando para añadir nota (invisible en el cuerpo)
    \newcommand{\eqnote}[1]{%
      \addcontentsline{leq}{leqnote}{#1}%
    }
    
    % Comando para ecuaciones
    \newcommand{\eqblock}[2]{%
      % Verificar si necesitamos imprimir el título de sección
      \ifx\leq@current@section\leq@last@printed@section\else
        \addcontentsline{leq}{leqsection}{\leq@current@section}%
        \global\let\leq@last@printed@section\leq@current@section
        \gdef\leq@last@printed@subsection{}% Reset subsección al cambiar sección
      \fi
      % Verificar si necesitamos imprimir el título de subsección
      \ifx\leq@current@subsection\leq@last@printed@subsection\else
        \ifx\leq@current@subsection\@empty\else
          \addcontentsline{leq}{leqsubsection}{\leq@current@subsection}%
          \global\let\leq@last@printed@subsection\leq@current@subsection
        \fi
      \fi
      % Ecuación
      \refstepcounter{eqnum}%
      \begin{equation*}
        #1\tag{E.\theeqnum}
      \end{equation*}%
      \addcontentsline{leq}{leqt}{[E.\theeqnum] -- #2}%
      \addcontentsline{leq}{leqm}{#1}%
    }
    
    %%% ----- Índice de ecuaciones ----------%%%
    % Tamaño para []
    \newcommand{\leqIndexFont}{\normalsize}
    \newcommand{\leqTitleFont}{\tiny}
    \newcommand{\leqNoteFont}{\tiny}
    % Cabecera de sección 
    \newcommand*\l@leqsection[2]{%
      \addvspace{25pt}% Mayor separación antes de nueva sección
      \noindent\leqIndexFont
      \makebox[\linewidth]{\rule{.38\linewidth}{0.4pt}\ \ \textbf{  \MakeUppercase{#1}}\ \ \rule{.38\linewidth}{0.4pt}}%
      \par\vspace{-10pt}
    }
    % Cabecera de subsección 
    \newcommand*\l@leqsubsection[2]{%
      \par\addvspace{15pt}%
      \noindent\leqIndexFont #1
      \par\vspace{-7pt}
    }
    % Título de ecuación 
    \newcommand*\l@leqt[2]{%
    \par\addvspace{10pt}%
      \par\noindent\leqTitleFont #1\par
    }
    % Miniatura de ecuación
    \newcommand*\l@leqm[2]{%
      \vspace{0.3\baselineskip}%
      \par\scriptsize\noindent\hspace*{1.5em}\(#1\)\par%
    }
    % Nota de ecuación 
    \newcommand*\l@leqnote[2]{%
      \vspace{0.2\baselineskip}%
      \par\noindent\leqNoteFont\hspace*{1.5em}\textbf{Nota.--} #1\par\vspace{0.5\baselineskip}%
    }
    % Generación del índice
    \newcommand{\mylistofequations}{%
      \clearpage
      \phantomsection
      % Título visual de la lista (ajústalo a tu gusto):
      {\centering\bfseries\Large Índice de ecuaciones\par}
      % Entrada en nuestro índice 'stc':
      \addcontentsline{stc}{section}{Índice de ecuaciones}%
      % Llamada a tu lista real (usa el nombre exacto de tu macro):
      \@starttoc{leq}%
    }
    
    %--------- Captura de títulos de sección --------%
    \let\leq@original@section\section
    \let\leq@original@subsection\subsection
    % Flag para saber si estamos en una sección asterisco
    \newif\ifLeqStarSection
    \LeqStarSectionfalse
    
    % Redefinir \section CON CONTROL DE ASTERISCO
    \renewcommand{\section}{%
      \@ifstar{\leq@section@star}{\@dblarg\leq@section@nostar}%
    }
    \def\leq@section@star#1{%
      \global\LeqStarSectiontrue
      \gdef\leq@current@section{#1}%
      \gdef\leq@current@subsection{}%
      \leq@original@section*{#1}%
      \global\LeqStarSectionfalse
    }
    \def\leq@section@nostar[#1]#2{%
      \global\LeqStarSectionfalse
      \gdef\leq@current@section{#2}%
      \gdef\leq@current@subsection{}%
      \leq@original@section[#1]{#2}%
    }
    % Redefinir \subsection
    \renewcommand{\subsection}{%
      \@ifstar{\leq@subsection@star}{\@dblarg\leq@subsection@nostar}%
    }
    \def\leq@subsection@star#1{%
      \global\LeqStarSectiontrue
      \gdef\leq@current@subsection{#1}%
      \leq@original@subsection*{#1}%
      \global\LeqStarSectionfalse
    }
    \def\leq@subsection@nostar[#1]#2{%
      \global\LeqStarSectionfalse
      \gdef\leq@current@subsection{#2}%
      \leq@original@subsection[#1]{#2}%
    }

    % ----- Restauración de valores Globales de estilo ----------%
    % Flags
    \newif\ifInFootnote
    \InFootnotefalse
    \pretocmd{\@footnotetext}{\InFootnotetrue}{}{}
    \apptocmd{\@footnotetext}{\InFootnotefalse}{}{}
    % Profundidad de listas (soporta anidadas)
    \newcount\MyListDepth \MyListDepth=0
    \AtBeginEnvironment{la}{\advance\MyListDepth by 1}
    \AfterEndEnvironment{la}{\advance\MyListDepth by -1}
    \AtBeginEnvironment{lnum}{\advance\MyListDepth by 1}
    \AfterEndEnvironment{lnum}{\advance\MyListDepth by -1}
    \AtBeginEnvironment{itemize}{\advance\MyListDepth by 1}
    \AfterEndEnvironment{itemize}{\advance\MyListDepth by -1}
    \AtBeginEnvironment{enumerate}{\advance\MyListDepth by 1}
    \AfterEndEnvironment{enumerate}{\advance\MyListDepth by -1}
    % Restauración diferida: hazlo al INICIO del próximo párrafo real
    \newif\if@pendingRestore
    \@pendingRestorefalse
    \newcommand{\RequestRestore}{\global\@pendingRestoretrue}
    \everypar{%
      \if@pendingRestore
        \global\@pendingRestorefalse
        \ifInFootnote\else
          \ifnum\MyListDepth=0
            \setlength{\parskip}{\GlobalParskip}%
            \setstretch{\GlobalBaseLineStretch}%
          \fi
        \fi
      \fi
    }
    % Al final de bloques:
    \AfterEndEnvironment{equation}{\RequestRestore}
    \AfterEndEnvironment{equation*}{\RequestRestore}
    \AfterEndEnvironment{la}{\RequestRestore}
    \AfterEndEnvironment{lnum}{\RequestRestore}
    % Al final de macros:
    \apptocmd{\eqblock}{\RequestRestore}{}{}
    \apptocmd{\imagenfit}{\RequestRestore}{}{}
    
\makeatother

% ===================== ToC propio con 4 niveles (único bloque) =====================
\makeatletter

% --- Escritores: añadimos una línea al .stc en cada cabecera numerada ---
% Guardamos las versiones actuales para llamar al comportamiento normal
\let\MyTOC@orig@section\section
\let\MyTOC@orig@subsection\subsection
\let\MyTOC@orig@subsubsection\subsubsection
\let\MyTOC@orig@paragraph\paragraph

% SECTION
\renewcommand{\section}{%
  \@ifstar{\MyTOC@section@star}{\@dblarg\MyTOC@section@nostar}%
}
\def\MyTOC@section@star#1{%
  \MyTOC@orig@section*{#1}% sin entrada al .stc
}
\def\MyTOC@section@nostar[#1]#2{%
  \MyTOC@orig@section[{#1}]{#2}% comportamiento normal (escribe al .toc)
  % Escribimos también nuestra línea al .stc (texto con \numberline)
  \addcontentsline{stc}{section}{\protect\numberline{\thesection}#2}%
}

% SUBSECTION
\renewcommand{\subsection}{%
  \@ifstar{\MyTOC@subsection@star}{\@dblarg\MyTOC@subsection@nostar}%
}
\def\MyTOC@subsection@star#1{%
  \MyTOC@orig@subsection*{#1}%
}
\def\MyTOC@subsection@nostar[#1]#2{%
  \MyTOC@orig@subsection[{#1}]{#2}%
  \addcontentsline{stc}{subsection}{\protect\numberline{\thesubsection}#2}%
}

% SUBSUBSECTION
\renewcommand{\subsubsection}{%
  \@ifstar{\MyTOC@subsubsection@star}{\@dblarg\MyTOC@subsubsection@nostar}%
}
\def\MyTOC@subsubsection@star#1{%
  \MyTOC@orig@subsubsection*{#1}%
}
\def\MyTOC@subsubsection@nostar[#1]#2{%
  \MyTOC@orig@subsubsection[{#1}]{#2}%
  \addcontentsline{stc}{subsubsection}{\protect\numberline{\thesubsubsection}#2}%
}

% PARAGRAPH
\renewcommand{\paragraph}{%
  \@ifstar{\MyTOC@paragraph@star}{\@dblarg\MyTOC@paragraph@nostar}%
}
\def\MyTOC@paragraph@star#1{%
  \MyTOC@orig@paragraph*{#1}%
}
\def\MyTOC@paragraph@nostar[#1]#2{%
  \MyTOC@orig@paragraph[{#1}]{#2}%
  % Si no numeras los \paragraph, cambia \theparagraph por texto plano sin \numberline
  \addcontentsline{stc}{paragraph}{\protect\numberline{\theparagraph}#2}%
}

% --- Impresor del índice propio (lee el .stc) ---
\newcommand{\MyTOC}{%
  \par\bigskip
  {\centering\bfseries\Large Índice de contenido\par}%
  \medskip
  \begingroup
    \parindent=0pt\parskip=2pt
    \@starttoc{stc}%
  \endgroup
}

\makeatother
% ===================== fin ToC propio =====================


%%%%%%%%%%%%%%


% --- Datos del tema ---
\title{Sistema Fiscal en España (IV): Imposición Patrimonial\\
\large El Impuesto sobre el Patrimonio. El Impuesto sobre Sucesiones y Donaciones. Impuesto sobre Bienes Inmuebles. Impuesto sobre Transmisiones Patrimoniales y Actos Jurídicos Documentados}
\author{Marcelo Ortega}
\date{Septiembre 2026}

\begin{document}
\maketitle
\newpage

% --- Página de resumen y modificaciones ---
\puntosclave{ %% Escribir puntos clave Apuntes CECO
     
    }
\vspace{1cm}

\modificaciones{
    
    }

\newpage
\MyTOC
\newpage

\mylistofequations
\clearpage

\MyListOfFigures
\clearpage


\pagenumbering{arabic}
\setcounter{page}{1}


\section{Introducción}



\subsection{Contextualización}


\subsubsection{Relevancia}




\subsubsection{Problemática}



\subsection{Estructura de la exposición}







\clearpage
\section{Imposición patrimonial: Marco conceptual}
\puntosclave{
    
}   

Cuatro preguntas vertebran este apartado, en el orden en que conviene responderlas: de dónde viene esta forma de imposición (1.1); por qué existe, con el debate clásico y el debate de máxima actualidad desarrollados juntos, sin fragmentar (1.2); cómo se articula en la práctica española, presentando brevemente las cuatro figuras y el criterio que las diferencia, para detenerme en el problema de fondo que esa articulación genera —la superposición de varias figuras sobre un mismo patrimonio— (1.3); y, por último, qué ha ocurrido cuando ese marco formal único se ha fragmentado en diecisiete marcos materiales distintos, con el conflicto sociopolítico y de competencia fiscal que eso ha generado (1.4). Empiezo por la historia.

\subsection{Historia: ¿cómo se ha desarrollado la imposición patrimonial en España?}
\textcolor{red}{¿Un poco corto?}

La imposición patrimonial española no tiene una única genealogía, porque nunca fue concebida como una sola figura: la línea de la tenencia y la línea de la transmisión evolucionan por separado durante la mayor parte de su historia, y solo convergen, de forma tardía, en el diseño competencial que hoy conocemos.

[Seguro] La línea de la tenencia sobre bienes inmuebles es la más antigua de las dos: se remonta a la Contribución Territorial, rústica y urbana, que ya mencionamos en el Tema 16 como una de las figuras del viejo sistema de imposición de producto anterior a 1900, y que sigue existiendo, transformada, como antecedente directo del actual IBI —una continuidad de más de un siglo que contrasta con la juventud relativa del resto de figuras de este tema—. La tenencia de patrimonio en sentido general, en cambio, no tuvo impuesto propio hasta muy avanzado el siglo XX: fue la Ley 50/1977, de 14 de noviembre —la misma norma que ya situamos en el Tema 17 como antesala de la reforma societaria de 1978— la que creó el Impuesto Extraordinario sobre el Patrimonio de las Personas Físicas, concebido en su origen, como su propio nombre indica, como figura excepcional y censal: su función declarada no era tanto recaudatoria como de control, permitir a la Hacienda conocer con precisión el patrimonio de los contribuyentes para contrastarlo con la renta recién sometida al IRPF de 1978, en un contexto de fraude fiscal masivo y de ausencia casi total de información tributaria fiable sobre las grandes fortunas españolas. [Seguro] El carácter "extraordinario" se abandonó formalmente con la Ley 19/1991, de 6 de junio, que lo convirtió en impuesto ordinario y permanente sobre el patrimonio neto, coincidiendo cronológicamente —no por azar, sino como parte del mismo paquete legislativo— con la reforma general del IRPF de ese mismo año ya analizada en el Tema 16. [Seguro] Su recorrido posterior muestra una volatilidad política poco habitual en el sistema tributario español: la Ley 4/2008 lo suprimió de facto mediante una bonificación del 100% de la cuota, manteniéndolo formalmente vigente pero sin efecto recaudatorio durante los ejercicios 2008 a 2010; el Real Decreto-ley 13/2011 lo reinstauró, con carácter inicialmente temporal y sucesivas prórrogas anuales que se han convertido, de facto, en permanencia indefinida, en el contexto de necesidad recaudatoria de la crisis de deuda soberana; y el Real Decreto-ley 3/2022 creó el Impuesto Temporal de Solidaridad de las Grandes Fortunas como parche estatal armonizador —figura que desarrollaremos con detalle en el apartado 2, y cuya sola existencia ya anuncia que algo se había roto en el diseño territorial del IP, aunque el porqué de esa ruptura lo reservo, con toda su importancia, para el subapartado 1.4—.

[Seguro] La línea de la transmisión tiene un origen distinto y considerablemente más antiguo en su núcleo conceptual: desciende del histórico Impuesto de Derechos Reales y Transmisión de Bienes, documentado ya en reformas de la Hacienda española de principios del siglo XX —una reforma de 1926, por ejemplo, introdujo dentro de esa misma figura el gravamen sobre el caudal relicto indiviso de la herencia, anticipando ya entonces el debate entre gravar el caudal global o las porciones hereditarias individuales que retomaremos con profundidad en el apartado 3—. No desarrollo más este antecedente aquí, porque es precisamente el puente que reservo para la introducción del apartado 3: baste fijar que ISD e ITPAJD no son dos impuestos que coincidan por azar en gravar transmisiones, sino que comparten un tronco normativo común del que solo se escindieron de forma definitiva con la Ley 32/1980 (ITPAJD) y la Ley 29/1987 (ISD), ya en pleno periodo constitucional.

[Seguro] El tramo final de esta historia —que aquí solo dejo apuntado, porque merece desarrollo propio en el subapartado 1.4— es el de la cesión progresiva a las Comunidades Autónomas, que afecta a IP, ISD e ITPAJD (no al IBI, que nunca dejó de ser local) y que se articula en escalones normativos sucesivos de alcance creciente: la Ley Orgánica 8/1980 (LOFCA) sienta el marco general de los tributos cedidos, pero la cesión efectiva durante toda la década de los ochenta y buena parte de los noventa fue, en la práctica, solo de rendimiento y de gestión; el salto a una capacidad normativa real llega con la Ley 14/1996, se amplía sustancialmente con la Ley 21/2001 —que permite a las Comunidades reducir de forma sustancial el ISD, posibilidad que unas aceptaron y otras rechazaron, sembrando ya entonces la semilla de la dispersión que hoy conocemos— y culmina, en su diseño vigente, con la Ley 22/2009. Es en este punto donde la historia de la imposición patrimonial deja de ser una historia de normas estatales sucesivas y se convierte en diecisiete historias paralelas — el objeto exacto del subapartado 1.4.



\subsection{¿Por qué gravar el patrimonio?}
## 1.2. ¿Por qué? Justificación teórica de la imposición patrimonial

Si la renta ya tributa en el IRPF y en el IS, ¿qué añade gravar, además, el patrimonio que esa renta ha ido generando y acumulando? Esta pregunta no es retórica ni exclusivamente española: es el debate de fondo que explica por qué la inmensa mayoría de los países de la OCDE que tuvieron alguna vez un impuesto general sobre el patrimonio lo han ido abandonando desde los años noventa —Alemania en 1997, Austria y Dinamarca antes, Francia lo sustituyó en 2018 por un impuesto limitado a inmuebles (IFI)—, mientras España, Noruega, Suiza y algún otro país lo mantienen en su forma general. Conviene recorrer los argumentos a favor y en contra con el mismo detenimiento que la propia doctrina les ha dedicado, porque la respuesta no es evidente ni unánime.

[Seguro] El argumento clásico más citado, ya lo anticipamos en el Tema 16 al hablar de Henry Simons a propósito de la renta, es el de la **capacidad económica autónoma**: la titularidad de un patrimonio genera, por sí misma, un poder de disposición, una seguridad económica y una capacidad de acceso al crédito que no dependen de que ese patrimonio genere o no un flujo de renta observable en un ejercicio concreto. Dos personas con la misma renta anual, pero con patrimonios acumulados radicalmente distintos, no tienen la misma capacidad económica real, aunque el IRPF las trate de forma idéntica: quien tiene patrimonio puede afrontar imprevistos, acceder a financiación en condiciones más ventajosas, y disfruta de una tranquilidad patrimonial que el mero flujo de renta no capta. [Probable] Emparentado con este argumento, pero distinto en su mecánica, está el modelo doctrinal de origen alemán que trata el impuesto sobre el patrimonio como una **Ergänzungssteuer** —impuesto complementario o de cierre—, cuya función no es gravar una capacidad económica nueva y distinta de la renta, sino **corregir por la puerta de atrás** las propias lagunas del impuesto sobre la renta que ya identificamos con detalle en el Tema 16: la renta imputada de la vivienda habitual que el IRPF no grava, las plusvalías latentes no realizadas que el criterio de realización deja fuera año tras año, el patrimonio heredado por la vía de donaciones fraccionadas que eluden parcialmente el ISD. Bajo esta lectura, el IP no es un impuesto verdaderamente autónomo, sino un mecanismo de cierre que grava, de forma indirecta y aproximada, exactamente las mismas divergencias estructurales del IRPF respecto de SHS que ya analizamos —una conexión que conviene que tengas presente, porque es la que explica mejor por qué un contribuyente puede percibir el IP como una "segunda vuelta" sobre una renta que el IRPF ya debería haber gravado y no gravó del todo—.

[Seguro] Un tercer argumento, de naturaleza más administrativa que sustantiva, es el que ya vimos operar históricamente en España: la **función de control**. El IP nació en 1977 no tanto para recaudar como para obligar al contribuyente a declarar, año tras año, el conjunto de su patrimonio, generando así un dato cruzado que la Administración puede contrastar con la renta declarada en el IRPF y con las adquisiciones patrimoniales declaradas en el ISD y el ITPAJD —una lógica de "trazabilidad" de la riqueza que sigue siendo, incluso hoy, uno de los argumentos que la doctrina hacendística esgrime a favor de mantener el impuesto, con independencia de su rendimiento recaudatorio directo—. Un cuarto argumento, de raíz más filosófico-política que económica, es el **igualitario**: la concentración patrimonial heredada de generación en generación perpetúa desigualdades de punto de partida que ni siquiera un sistema de renta perfectamente progresivo puede corregir por sí solo, porque la renta del trabajo de cada generación parte de una base patrimonial ya desigual —un argumento de raíz rawlsiana sobre la igualdad de oportunidades, distinto del argumento más específico de Piketty que desarrollo más abajo—. Y, en sentido inverso, cabe también un argumento de **teoría del beneficio**, análogo al que empleamos para justificar el IS en el Tema 17: quien acumula un patrimonio se beneficia de forma particularmente intensa de la protección jurídica que el Estado dispensa a la propiedad —el sistema judicial que hace valer sus derechos reales, el sistema financiero regulado que preserva el valor de sus activos, la seguridad pública que protege sus bienes—, lo que justificaría una contribución adicional y específica más allá de la que ya aporta como contribuyente del IRPF.

[Seguro] Frente a este conjunto de argumentos, la crítica clásica se articula en varios frentes bien documentados. El primero es la **doble imposición del ahorro**, la misma tensión que ya analizamos con Kaldor en el Tema 16: si la renta que dio origen al patrimonio ya tributó en su día por el IRPF, gravar después el patrimonio acumulado penaliza doblemente la decisión de ahorrar frente a la de consumir, exactamente el mismo problema de neutralidad intertemporal que allí desarrollamos, ahora proyectado sobre el stock y no sobre el flujo. El segundo es el problema de la **valoración e iliquidez**, especialmente agudo para activos empresariales o inmobiliarios sin mercado organizado: un empresario cuyo patrimonio está mayoritariamente inmovilizado en su empresa familiar puede verse obligado a generar liquidez artificial —vendiendo participaciones, endeudándose— solo para poder pagar un impuesto sobre un valor que, en la práctica, no puede materializar sin deshacerse de la propia fuente de su actividad, el mismo problema de restricción de liquidez que ya vimos, en otro contexto, a propósito del criterio de realización de las ganancias patrimoniales en el Tema 16. El tercer argumento, más técnico y menos conocido fuera de la literatura especializada pero de gran peso analítico, es el de la **tasa implícita sobre el rendimiento real**: un impuesto sobre el patrimonio del 1% anual puede parecer modesto en términos nominales, pero si el activo gravado rinde, por ejemplo, un 2% real anual, ese 1% equivale, en términos económicos, a un tipo efectivo del **50% sobre el rendimiento real** que el activo genera —una carga muy superior a la que sugiere el tipo nominal, y que se agrava en periodos de tipos de interés bajos o de rentabilidades reducidas, penalizando desproporcionadamente a quien tiene patrimonio pero bajo rendimiento frente a quien tiene menos patrimonio pero rentabilidades altas—. El cuarto es el ya mencionado argumento de **movilidad internacional del capital**: la abrogación alemana de 1997 y la sustitución francesa de 2018 se justificaron, en ambos casos, invocando la fuga de grandes patrimonios hacia jurisdicciones sin imposición patrimonial, el mismo mecanismo de competencia fiscal que ya analizamos con Zodrow y Mieszkowski en el Tema 17 a propósito de la dispersión foral, ahora proyectado a escala internacional.

[Seguro] El debate contemporáneo, el que le da a esta cuestión su actualidad más intensa, tiene dos nombres propios ineludibles. **Thomas Piketty**, en *Le Capital au XXIe siècle* (2013), documenta empíricamente que la rentabilidad del capital ha superado de forma sistemática, a lo largo de la historia económica reciente, la tasa de crecimiento de la economía —su célebre desigualdad **r > g**—, con la consecuencia de que la riqueza acumulada tiende a crecer más deprisa que la renta del trabajo, generando una concentración patrimonial creciente que ni el crecimiento económico ni la meritocracia laboral pueden compensar por sí solos. Su propuesta normativa —un impuesto progresivo sobre el capital de alcance idealmente global, para evitar precisamente la fuga que acabamos de mencionar— es el argumento de autoridad más citado, en ambos sentidos del debate, de la última década.

[Seguro] **Gabriel Zucman**, coautor habitual de Piketty y discípulo directo de esa misma escuela, ha desplazado el argumento del terreno teórico al empírico y, más recientemente, al institucional. En *The Hidden Wealth of Nations* (2015) documenta, con metodología propia de reconciliación de estadísticas de activos y pasivos entre países, la magnitud de la riqueza oculta en jurisdicciones offshore —un patrimonio que escapa sistemáticamente tanto al IRPF como a cualquier impuesto patrimonial nacional—, proporcionando el fundamento empírico de por qué cualquier imposición patrimonial puramente nacional está condenada a una elusión estructural si no se coordina internacionalmente. [Seguro] Esa misma lógica culmina, en junio de 2024, en el informe que Zucman elaboró por encargo de la **presidencia brasileña del G20**, titulado *A Blueprint for a Coordinated Minimum Effective Taxation Standard for Ultra-High-Net-Worth Individuals*: propone un **tipo mínimo efectivo del 2%** sobre el patrimonio neto de los aproximadamente 3.000 multimillonarios del planeta (patrimonio superior a 1.000 millones de dólares), calculando que estos soportan hoy un tipo efectivo medio de apenas el 0,3% de su riqueza —muy por debajo del que soporta cualquier asalariado medio sobre su renta—, con un potencial recaudatorio estimado entre 200.000 y 250.000 millones de dólares anuales. [Probable] El diseño técnico de la propuesta es deliberadamente **paralelo al Pilar 2 de la OCDE/G20 que ya analizamos en el Tema 17 para el Impuesto sobre Sociedades**: no crea un nuevo impuesto mundial sobre el patrimonio, sino un mecanismo de **complemento** (*top-up*) que se activa únicamente cuando la tributación efectiva combinada de renta y patrimonio de un multimillonario cae por debajo de ese estándar mínimo del 2% —el mismo diseño de "suelo mínimo coordinado" que vimos operar sobre los grupos multinacionales, ahora trasladado a las personas físicas de mayor patrimonio—. La propuesta ha recibido el apoyo expreso de Brasil, Francia, España y Sudáfrica dentro del G20, y la oposición declarada de Estados Unidos, en un debate que se prolonga sin resolución definitiva desde la cumbre de Río de Janeiro de noviembre de 2024.

[Seguro] El episodio que mejor ilustra, con nombres y cifras concretas, la naturaleza genuinamente controvertida —y no solo políticamente polarizada— de este debate es el enfrentamiento académico entre **Saez y Zucman**, de un lado, y **Lawrence Summers y Natasha Sarin**, del otro, a propósito de la propuesta de impuesto sobre el patrimonio de la senadora Elizabeth Warren durante las primarias demócratas de 2019. Saez y Zucman, autores materiales de la propuesta, estimaron que un impuesto del 2% sobre patrimonios superiores a 50 millones de dólares (con un tramo adicional del 1% por encima de 1.000 millones) recaudaría **212.000 millones de dólares anuales**, asumiendo, sobre la base de cuatro estudios académicos publicados sobre la experiencia de impuestos patrimoniales en otros países, una tasa de elusión del 15% del patrimonio afectado. Summers y Sarin, en sendos artículos de opinión en *The Washington Post* (abril y junio de 2019), cuestionaron esa cifra invocando la experiencia del propio impuesto de sucesiones estadounidense —que recauda históricamente muy por debajo de lo que su base teórica permitiría prever, precisamente por la magnitud de la planificación fiscal y la elusión que genera cualquier impuesto sobre grandes patrimonios— y llegaron a plantear que la recaudación real podría situarse en apenas 25.000 millones de dólares, una fracción de la estimación original. Saez y Zucman replicaron con dureza, acusando a sus críticos de partir de un supuesto implícito de elusión cercano al 90%, y el debate técnico —sobre qué tasa de elusión es razonable asumir a falta de datos administrativos reales sobre un impuesto que nunca ha existido en Estados Unidos— quedó, en lo sustancial, sin resolver, precisamente porque **la pregunta empírica central no tiene respuesta hasta que el impuesto se aplica de verdad**, un problema epistémico que afecta a cualquier debate sobre imposición patrimonial y que conviene que sepas trasladar al tribunal como la razón última de que este debate no se cierre nunca del todo con argumentos puramente teóricos.

[Probable] Cierro el subapartado devolviendo el argumento a España: el hecho de que nuestro país mantenga un impuesto general sobre el patrimonio, frente a la tendencia derogatoria mayoritaria en la OCDE, no responde a que España haya "ganado" este debate doctrinal de forma concluyente —como acabas de ver, no existe tal cosa—, sino a una combinación de necesidad recaudatoria coyuntural (su reinstauración en 2011 fue una respuesta a la crisis de deuda soberana, no una decisión de diseño tributario sereno) y de una arquitectura competencial que, como desarrollaremos con toda su importancia en el subapartado 1.4, ha terminado generando un problema añadido que ninguno de los argumentos doctrinales que acabas de leer anticipaba: la posibilidad de que una sola Comunidad Autónoma neutralice, por vía de bonificación, la decisión estatal de mantener el impuesto.


[Seguro] Tienes razón, y el motivo técnico es el mismo que ya fijamos para "dispersión foral" en el Tema 17: "clásica" y "actual" no son dos etapas de un único argumento que se construye progresivamente —donde partir el texto sería cortar una frase a la mitad—, sino **dos vertientes materialmente distintas** de la misma idea general ("por qué existe la imposición patrimonial"), cada una con su propio aparato de autores y su propio marco temporal. Ese es exactamente el criterio que debo aplicar de forma sistemática: subdividir en sub-subapartados cuando las facetas son **paralelas y enumerables** (aunque dialoguen entre sí), y mantener el desarrollo continuo cuando una idea se construye como una única cadena argumental en la que cortar generaría un corte artificial. Lo aplico ya en este subapartado.

\subsection{Materialización: ¿cómo se articula la imposición patrimonial en España?}

### 1.3.1. Las cuatro figuras y el criterio que las diferencia

Recuperando, sin volver a desarrollarlos, los dos ejes fijados al proponer la estructura del tema, España articula la imposición patrimonial en cuatro figuras normativas —IP (con su parche armonizador, el ITSGF), IBI, ISD e ITPAJD—, que se distribuyen así: por el eje del **momento del gravamen**, IP e IBI gravan la tenencia periódica (apartado 2), ISD e ITPAJD gravan la transmisión puntual (apartado 3); por el eje de **ámbito y causa**, dentro de la tenencia el IP es general y personal —todo el patrimonio neto del contribuyente— mientras el IBI es real y objetivo —solo bienes inmuebles—, y dentro de la transmisión el ISD grava la causa gratuita —herencias y donaciones— mientras el ITPAJD grava, en su modalidad TPO, la causa onerosa entre particulares, con AJD gravando la mera formalización documental con independencia de la causa. [Probable] Este doble criterio no es solo un instrumento clasificatorio previo a la redacción: es también, como vas a ver ahora, la clave para entender **por qué** estas cuatro figuras, pese a operar formalmente con total autonomía —cada una con su propio hecho imponible, su propia ley, en algunos casos su propia Administración competente—, terminan incidiendo, con enorme frecuencia, sobre el mismo bien y el mismo contribuyente casi de forma simultánea.

### 1.3.2. El problema de la superposición

La superposición dentro de esta familia de impuestos tiene varias manifestaciones, de intensidad y de tratamiento jurídico muy distintos, y conviene que sepas distinguirlas antes de que profundicemos en la que más lo merece.

Existe una **superposición horizontal dentro de la tenencia**: un mismo inmueble tributa, cada año, tanto por IBI —en su condición de bien real, con independencia de quién sea su titular— como, si ese titular es persona física y supera el mínimo exento, por IP —integrado en su patrimonio neto general—. Es una superposición de baja intensidad conflictiva, porque ambas figuras operan sobre magnitudes y en niveles de la Administración distintos (local frente a estatal-cedido), y su tratamiento detallado corresponde, como ya acordamos, al apartado 2. Existe también una **continuidad entre transmisión y tenencia futura**: el heredero que paga ISD por adquirir un inmueble empieza, desde ese mismo momento, a tributar por IP e IBI sobre ese mismo bien —no son impuestos simultáneos sobre el mismo hecho, sino una secuencia en la que el impuesto de transmisión abre la puerta a los de tenencia—, cuestión que retomaremos, brevemente, al cerrar el apartado 3.

[Seguro] Pero la superposición que merece desarrollo propio aquí, porque es la única que el legislador ha considerado tan grave como para necesitar una **corrección normativa expresa**, es la que se produce entre la imposición sobre la **renta** y la imposición sobre el **patrimonio**: el llamado **límite conjunto IRPF-IP**, regulado en el art. 31.Uno de la Ley 19/1991. Es la materialización jurídica exacta de la crítica de doble imposición del ahorro que desarrollamos en el subapartado 1.2 a propósito de Kaldor: si el mismo patrimonio ya generó renta gravada por el IRPF, y ese mismo patrimonio vuelve a tributar íntegramente por el IP, la suma de ambas cuotas podría llegar a absorber, en casos extremos, una proporción confiscatoria de la renta anual del contribuyente —imagina un patrimonio cuantioso pero de baja rentabilidad, invertido por ejemplo en activos que generan poca renta corriente: el titular podría verse obligado a liquidar parte de su patrimonio solo para poder pagar ambos impuestos con la renta que efectivamente percibe—. [Seguro] La norma responde fijando que la cuota íntegra del IP, sumada a las cuotas del IRPF (estatal y autonómica), no puede superar el **60% de la base imponible del IRPF** del contribuyente sujeto por obligación personal; si se supera ese límite, se reduce la cuota del IP hasta alcanzarlo, pero con un **suelo de protección recaudatoria**: esa reducción nunca puede superar el 80% de la cuota del IP, de modo que, como mínimo, el contribuyente sigue pagando el 20% de la cuota que le correspondería sin límite alguno. Es, otra vez, el mismo patrón que ya identificamos repetidas veces en los Temas 16 y 17: un principio estructural correctamente reconocido por el legislador, inmediatamente acotado por una salvaguarda que protege la recaudación frente a la aplicación plena de ese mismo principio.

[Seguro] Este mecanismo, lejos de ser una cuestión técnica cerrada, es hoy objeto de un doble frente de litigiosidad de gran actualidad que conviene que conozcas. El primero es de **Derecho comparado y de Derechos Humanos**: la doctrina española más reciente sobre la materia ha señalado que el Tribunal Europeo de Derechos Humanos ha entendido, en su jurisprudencia sobre imposición patrimonial en otros Estados, que una carga tributaria conjunta que supere el **50%** de la renta del contribuyente puede resultar confiscatoria y, por tanto, contraria al art. 1 del Protocolo Adicional 1 al Convenio Europeo de Derechos Humanos —lo que plantea la pregunta, todavía no resuelta con autoridad en España, de si el límite español del 60% (con su suelo del 20% que en la práctica puede empujar la carga conjunta bastante por encima de ese umbral) es compatible con esa jurisprudencia europea, sin que exista hasta la fecha un pronunciamiento del Tribunal Constitucional español que zanje la cuestión con la claridad que el problema merece—. [Seguro] El segundo frente es de **Derecho de la Unión Europea**, y es de rabiosa actualidad: el Tribunal Supremo, en sendas Sentencias de **29 de octubre y 3 de noviembre de 2025**, ha reconocido que los contribuyentes **no residentes** del IP —sujetos por obligación real, solo sobre los bienes situados en España— tienen también derecho a aplicar este límite conjunto, frente al criterio tradicional de la Administración tributaria, que lo reservaba exclusivamente a los residentes. El Tribunal Supremo funda esta extensión en el principio de libre circulación de capitales del art. 63 del Tratado de Funcionamiento de la Unión Europea, entendiendo que excluir a los no residentes de una salvaguarda anti-confiscatoria de la que sí disfrutan los residentes constituye una discriminación por razón de residencia no justificada dentro del mercado interior europeo —el mismo tipo de argumento de no discriminación transfronteriza que ya vimos operar, en otro contexto, en el Tema 17 a propósito de las deducciones por doble imposición internacional del IS—.

[Probable] Este subapartado cierra, por tanto, con una idea que conviene que traslades con claridad al tribunal: la superposición dentro de la imposición patrimonial no es un defecto de diseño que el legislador haya ignorado, sino un problema que el propio sistema reconoce y trata de corregir mediante mecanismos expresos —el límite conjunto es la prueba de que el legislador español sí tomó en serio la crítica doctrinal de doble imposición—, pero esa corrección genera, a su vez, nuevos frentes de litigio —el umbral exacto, su alcance subjetivo— que no estaban resueltos cuando la norma se aprobó en 1991 y que los tribunales, tanto españoles como europeos, siguen perfilando treinta años después.


\subsection{La descentralización: un marco formal, diecisiete marcos materiales}

### 1.4.1. La cesión y sus consecuencias: qué pueden hacer realmente las Comunidades Autónomas

Ya dejamos apuntada, al cerrar el desarrollo histórico, la secuencia normativa de la cesión —LOFCA 1980, Ley 14/1996, Ley 21/2001, Ley 22/2009—; toca ahora dar contenido sustantivo a esa secuencia, porque el dato relevante no es que existan tres escalones, sino **qué capacidad normativa real** otorga cada uno sobre cada una de las tres figuras cedidas. [Probable] En el IP, las Comunidades Autónomas pueden regular el mínimo exento, la tarifa y las deducciones y bonificaciones de la cuota —esta última competencia es la que ha resultado decisiva en la práctica—: **Madrid** aplica, desde 2008, una bonificación del 100% de la cuota, que neutraliza por completo el impuesto para cualquier contribuyente residente en su territorio con independencia de la cuantía de su patrimonio; **Andalucía** aprobó una bonificación equivalente desde 2022. En el ISD, la capacidad normativa alcanza reducciones propias en la base imponible, la tarifa, los coeficientes multiplicadores por patrimonio preexistente, y deducciones y bonificaciones en cuota, lo que ha generado —como ya vimos con detalle en el Tema 16— un mapa en el que la misma sucesión entre padres e hijos puede tributar de forma prácticamente simbólica en unas Comunidades y de forma sustancial en otras. En el ITPAJD, la capacidad se centra en los tipos de gravamen de TPO y AJD, con una dispersión de tipos —especialmente relevante en las transmisiones inmobiliarias, el hecho imponible de mayor volumen recaudatorio de la figura— que oscila, según la Comunidad, en varios puntos porcentuales sobre el mismo tipo de operación.

### 1.4.2. El debate sociopolítico: ¿autonomía fiscal legítima o competencia desleal interna?

[Probable] Este mapa no es un accidente técnico, sino la expresión de dos posiciones políticas e ideológicas enfrentadas sobre la propia función de estos impuestos, que conviene presentar con el mismo rigor con que presentamos el debate Tiebout/Zodrow-Mieszkowski en el Tema 17 a propósito de la competencia fiscal foral en el IS —el mismo argumento de fondo, trasladado ahora a la competencia entre Comunidades de régimen común—. Quienes defienden las bonificaciones —articulado con particular intensidad en el debate público en torno al ISD, donde ha calado la expresión, importada del debate anglosajón, del **"impuesto de la muerte"**— sostienen que gravar la transmisión mortis causa penaliza el ahorro familiar ya tributado, que resulta especialmente gravoso para el patrimonio inmobiliario ilíquido de clase media (la vivienda familiar heredada, sin capacidad de generar liquidez para pagar el impuesto), y que la competencia entre Comunidades por reducirlo es, en sí misma, una forma legítima y deseable de disciplina fiscal sobre gobiernos regionales que de otro modo tenderían a gastar en exceso —el argumento clásico de la *public choice* aplicado al federalismo fiscal—. Quienes lo critican sostienen que esta competencia no responde a un ejercicio genuino de autonomía política sobre el **nivel de gasto público** deseado por cada territorio —que sería la justificación clásica de la descentralización fiscal, y que corresponde desarrollar en el Tema 14—, sino a una carrera hacia el fondo sobre impuestos que gravan, casi en exclusiva, a las rentas y patrimonios más altos y más móviles, con el resultado de que la carga fiscal efectiva sobre la riqueza termina dependiendo, en última instancia, no de la capacidad económica del contribuyente sino de su lugar de residencia —una vulneración material, aunque no formal, del principio de igualdad del art. 31.1 CE que ya identificamos como tensión de fondo en el Tema 16 a propósito de la dispersión autonómica del IRPF—.

### 1.4.3. La respuesta estatal armonizadora: el ITSGF y su blindaje constitucional

[Seguro] La respuesta legislativa a este conflicto llegó con la **Ley 38/2022, de 27 de diciembre**, que en su art. 3 crea el **Impuesto Temporal de Solidaridad de las Grandes Fortunas**: un tributo estatal, complementario del IP y **no susceptible de cesión a las Comunidades Autónomas**, que grava con una cuota adicional los patrimonios netos superiores a 3.000.000 de euros. [Seguro] Su mecánica es la clave técnica de todo el conflicto: el contribuyente calcula su cuota del ITSGF conforme a una tarifa estatal única, pero puede **deducir de esa cuota lo efectivamente pagado por IP** en su Comunidad Autónoma de residencia. La consecuencia es que, allí donde la Comunidad ha bonificado el IP al 100% —Madrid, Andalucía—, el contribuyente no tiene nada que deducir, y termina pagando íntegramente la cuota del ITSGF; allí donde la Comunidad no bonifica, el contribuyente deduce lo ya pagado por IP y solo abona, si acaso, la diferencia hasta la cuota estatal. [Probable] El resultado neto es que la carga fiscal final tiende a **igualarse en todo el territorio de régimen común**, con independencia de la política de bonificaciones autonómica —lo que explica que las Comunidades recurrentes calificaran la norma, en sus propios escritos de recurso, como un impuesto **"clon"** del IP, diseñado específicamente para neutralizar, por la vía de una figura estatal no cedida, una competencia normativa —la bonificación del IP— que sí les pertenece legítimamente por cesión—.

[Seguro] Las Comunidades de Madrid, Andalucía, Murcia y Galicia interpusieron sendos recursos de inconstitucionalidad contra el art. 3 de la Ley 38/2022, fundados en cuatro motivos principales: vulneración de la autonomía financiera autonómica y del bloque de constitucionalidad en materia de tributos cedidos; infracción del derecho de representación política del art. 23.2 CE (por la forma de tramitación parlamentaria de la norma); vulneración del principio de lealtad institucional del art. 2.1.g) LOFCA; e infracción de la seguridad jurídica del art. 9.3 CE, a la que se sumaron, en el fondo del asunto, alegaciones de vulneración de la capacidad económica, confiscatoriedad y retroactividad. [Seguro] El **Tribunal Constitucional**, en su **Sentencia 149/2023, de 7 de noviembre** (resolviendo el recurso de Madrid, ponente la magistrada Balaguer Callejón), y en la **Sentencia 170/2023, de 22 de noviembre** (resolviendo el recurso de Andalucía, ponente el magistrado Campo Moreno), desestimó íntegramente ambos recursos, doctrina después aplicada a los recursos paralelos de Murcia y de la propia Asamblea de Madrid en diciembre de 2023. [Probable] El razonamiento central del Tribunal —y esto es lo que conviene que domines con precisión, porque es el argumento que zanja, al menos en el plano constitucional interno, todo el debate de este subapartado— es que el IP es un tributo **cedido**, pero no **transferido**: el Estado conserva en todo momento la titularidad del impuesto y la potestad originaria para establecer, con carácter general, tributos complementarios que gravan la misma capacidad económica, sin que la cesión de capacidad normativa sobre un impuesto existente limite la potestad del Estado para crear, ex novo, una figura estatal distinta —por más que su diseño técnico neutralice, en la práctica, el efecto de las bonificaciones autonómicas—.

[Probable] Conviene, no obstante, que cierres este subapartado sin dar el conflicto por completamente resuelto, porque no lo está: la propia doctrina especializada que analizó las sentencias de 2023 señala que quedan **abiertos frentes de discusión distintos del constitucional interno**, singularmente en el terreno del Derecho de la Unión Europea —posible discriminación de contribuyentes no residentes frente a residentes en la aplicación del impuesto, y una eventual vulneración del principio de seguridad jurídica de la Unión por la forma imprevista y retroactiva en que la norma se introdujo, aplicándose a todo el ejercicio 2022 pese a aprobarse el 27 de diciembre de ese mismo año, cuando el periodo impositivo estaba prácticamente agotado—. [Suposición] Es exactamente el mismo patrón de litigiosidad en dos niveles —constitucional interno primero, Derecho de la Unión después— que ya vimos, en el Tema 17, a propósito de los límites del decreto-ley en el IS y de la fiscalidad foral vasca: el Tribunal Constitucional español zanja una batalla, pero dentro de la Unión Europea puede quedar abierta otra distinta, con fundamento jurídico y argumentos completamente diferentes.




\clearpage
\section{Imposición sobre el Patrimonio: Tenencia}
\puntosclave{
    
}

Estas dos figuras —o tres, si contamos el ITSGF como pieza separada del IP— comparten el rasgo que ya fijamos en el subapartado 1.3: gravan la mera titularidad de patrimonio, año tras año, con independencia de que se produzca o no transmisión alguna. Pero conviene no perder de vista, desde el primer momento, que esa nota común esconde una diferencia de fondo entre ambas: el IP es un impuesto general y personal, que mide la capacidad económica global del contribuyente; el IBI es un impuesto real y objetivo, que grava un bien concreto con independencia de las circunstancias personales de su titular. Trato primero el IP, con el ITSGF como su complemento armonizador ya presentado en el apartado 1, y desarrollo después el IBI, cerrando con la superposición horizontal entre ambos que ya anticipamos.

\subsection{Impuesto sobre el Patrimonio (y el ITSGF)}
El IP se regula en la Ley 19/1991, de 6 de junio, ya citada en el apartado 1 a propósito de su origen y de su volatilidad reciente; el ITSGF, en el art. 3 de la Ley 38/2022, de 27 de diciembre, también ya presentado en el subapartado 1.4. 


[Seguro] El hecho imponible del IP es la titularidad, por el sujeto pasivo, de un patrimonio neto superior al mínimo exento a fecha de devengo —el 31 de diciembre de cada año, sin periodo impositivo propiamente dicho, a diferencia de los impuestos periódicos que hemos visto en temas anteriores—; el patrimonio neto se calcula como la diferencia entre el valor de los bienes y derechos de contenido económico de que sea titular el contribuyente y las cargas y deudas que disminuyan su valor. [Seguro] El mínimo exento general es de 700.000 euros, aunque las Comunidades Autónomas pueden elevarlo —ejerciendo, en sentido contrario a la bonificación total, la misma capacidad normativa que ya analizamos en el subapartado 1.4—, y se suma a la exención de la vivienda habitual hasta 300.000 euros, que opera de forma independiente. [Probable] En cuanto a la valoración, los bienes inmuebles se valoran, desde la reforma antifraude de 2021 ya vista en el apartado 1, por el mayor de tres valores: el catastral, el comprobado por la Administración a efectos de otros tributos, o el de adquisición —con el valor de referencia catastral operando, en la práctica, como umbral de referencia dominante para los inmuebles adquiridos con posterioridad a la reforma—; los valores mobiliarios cotizados se valoran a cotización media del último trimestre, y los no cotizados según reglas específicas de valor teórico contable, salvo que resulten exentos por la vía que desarrollo a continuación.

[Seguro] Conviene fijar, junto al mínimo exento ya señalado, dos elementos estructurales del impuesto que hasta ahora habíamos usado sin definir. El art. 5 LIP distingue entre obligación personal de contribuir —la de los residentes en España, que tributan por la totalidad de su patrimonio neto con independencia de dónde se sitúen los bienes o puedan ejercitarse los derechos— y obligación real —la de los no residentes, que tributan únicamente por los bienes y derechos de que sean titulares y que estén situados, puedan ejercitarse o hubieran de cumplirse en territorio español—, distinción que ya empleamos, sin desarrollarla, al analizar el "escudo fiscal" del art. 31 LIP y su reciente extensión jurisprudencial a los no residentes en el subapartado 1.3. La tarifa estatal (art. 30 LIP), aplicable en defecto de escala autonómica propia, se articula en ocho tramos progresivos que arrancan en un tipo del 0,2% para los primeros 167.129,45 euros de base liquidable y alcanzan, en el tramo superior, un tipo marginal del 3,5% —elevado desde el 2,5% anterior por la Ley de Presupuestos para 2023, elevación que ha sido, de hecho, objeto de un recurso de inconstitucionalidad promovido por la patronal catalana Foment del Treball y un grupo de diputados, por entender que ese incremento marginal podría vulnerar el principio de no confiscatoriedad—. [Probable] Conviene que sepas que esta tarifa estatal es, en la práctica, minoritaria: la mayoría de las Comunidades Autónomas que no bonifican el impuesto al 100% han aprobado escalas propias, generalmente más gravosas que la estatal en los tramos altos, ejerciendo la misma capacidad normativa que ya documentamos con detalle en el subapartado 1.4 para el conjunto del impuesto.

2.1.2. La exención de empresa familiar: diseño generoso y litigio constante

[Seguro] La exención más relevante del impuesto, tanto por su volumen económico como por su litigiosidad, es la de las participaciones en entidades familiares (art. 4.Ocho.Dos LIP), sujeta a tres requisitos acumulativos: que la entidad realice una actividad económica real y no se limite a la mera gestión de un patrimonio mobiliario o inmobiliario —con un test objetivo, "más de la mitad del activo son valores" durante más de 90 días del ejercicio, formalmente idéntico al test de "entidad patrimonial" que ya analizamos en el Tema 17 a propósito del IS, una convergencia conceptual entre ambos impuestos que rara vez se señala pero que revela que el legislador usa el mismo criterio técnico para dos finalidades distintas—; que el sujeto pasivo participe individualmente en, al menos, el 5% del capital, o conjuntamente con su grupo de parentesco —cónyuge, ascendientes, descendientes y colaterales de segundo grado, ni uno más— en, al menos, el 20%; y que el propio sujeto pasivo, o cualquier miembro de ese grupo de parentesco, ejerza efectivamente funciones de dirección en la entidad, percibiendo por ellas una remuneración que represente más del 50% de la totalidad de sus rendimientos empresariales, profesionales y del trabajo.

[Probable] Cada uno de estos tres requisitos ha generado una jurisprudencia y una doctrina administrativa muy extensa, prueba de que el diseño, pese a su apariencia objetiva, deja amplios márgenes de interpretación que la práctica se ha encargado de explotar. La Dirección General de Tributos y el Tribunal Supremo han confirmado, en consultas y sentencias reiteradas (entre ellas las SSTS 1776/2016 y 1198/2016), que basta con que las funciones de dirección las ejerza cualquier miembro del grupo de parentesco, sin necesidad de que sea el titular de las participaciones quien las ejerza personalmente —lo que permite, en la práctica, estructuras donde el patriarca o la matriarca de la familia mantiene la titularidad mayoritaria mientras un hijo, por ejemplo, ejerce la dirección efectiva y cobra por ello, cumpliendo así el requisito sin que quien realmente controla el capital tenga que trabajar en la empresa—. [Seguro] El límite de esta flexibilidad lo marca, con firmeza creciente en la doctrina administrativa más reciente, la delimitación estricta del grupo de parentesco: la propia DGT ha denegado la exención, en una consulta vinculante de 2025, en un supuesto en el que las funciones de dirección habían pasado de un tío a un sobrino, precisamente porque el parentesco colateral de tercer grado —sobrinos respecto de tíos— queda fuera del grupo de parentesco legalmente tasado, por más que la sustitución generacional responda, en términos empresariales reales, a la misma lógica de continuidad familiar que la norma pretende proteger. [Suposición] Este contraste —flexibilidad amplia dentro del grupo de parentesco tasado, rigidez estricta en sus límites exactos— es el mejor ejemplo de un problema de diseño más general que la doctrina especializada en planificación patrimonial ha señalado repetidamente: una exención concebida para proteger la continuidad de la empresa familiar genuina puede, con la ingeniería societaria adecuada, extenderse a estructuras de tenencia de patrimonio con muy poca sustancia empresarial real, mientras que, paradójicamente, puede negarse en supuestos de sucesión familiar genuina que no encajan exactamente en la definición legal de parentesco —un desajuste entre la finalidad declarada de la norma y su aplicación literal que ya vimos, con otro objeto, al hablar de la reserva de empresa familiar en el ISD en el Tema 16.

2.1.3. El debate doctrinal de fondo: ¿impuesto general sobre el patrimonio, o impuesto limitado a lo inmobiliario?

[Probable] Más allá del debate general sobre la existencia misma de la imposición patrimonial que desarrollamos en el subapartado 1.2, existe un debate de diseño más específico, y de gran actualidad comparada, sobre qué forma concreta debe adoptar un impuesto sobre el patrimonio si se decide mantenerlo: ¿debe gravar la totalidad del patrimonio del contribuyente —inmuebles, valores mobiliarios, participaciones empresariales, depósitos—, como hace España, o debe limitarse, como hace Francia desde 2018 con su Impôt sur la Fortune Immobilière (IFI), a los activos inmobiliarios, dejando fuera el patrimonio financiero y empresarial? [Probable] Los defensores de la fórmula limitada —que es, en la práctica, hacia donde ha derivado la mayoría de la literatura de reforma fiscal comparada en la última década— argumentan que el patrimonio inmobiliario es, de todos los activos, el menos móvil internacionalmente y el más fácilmente verificable por la Administración —no puede trasladarse a otra jurisdicción como sí puede hacerlo una cartera de valores o una cuenta bancaria—, de modo que limitar el impuesto a este tipo de activo minimiza el riesgo de fuga de capitales que ya vimos protagonizar la abolición alemana de 1997, sin renunciar del todo a gravar una manifestación evidente y muy visible de capacidad económica. [Probable] Los defensores del modelo general —posición que sigue siendo la dominante en la literatura más próxima a Piketty y Zucman, ya desarrollada en el subapartado 1.2— replican que limitar el impuesto a lo inmobiliario introduce una distorsión de neutralidad entre formas de patrimonio completamente ajena a cualquier principio de capacidad económica: dos contribuyentes con idéntico patrimonio neto, uno invertido en una vivienda y otro en una cartera de acciones de igual valor, tributarían de forma radicalmente distinta bajo un modelo tipo IFI, pese a tener exactamente la misma capacidad económica —el mismo tipo de inequidad horizontal por forma de tenencia que ya identificamos, en otro contexto, al hablar de la ruptura de la comprehensividad de fuente en el Tema 16—, y que, además, un impuesto limitado a lo inmobiliario termina gravando de forma desproporcionada a la clase media propietaria de vivienda, cuyo patrimonio está mayoritariamente inmovilizado en el inmueble familiar, mientras exonera a las grandes fortunas cuyo patrimonio está, con frecuencia, mayoritariamente diversificado en activos financieros. [Suposición] España, al mantener el modelo general con la exención de empresa familiar como válvula de escape específica para el patrimonio productivo, ha optado, de facto, por una tercera vía intermedia entre ambos extremos: grava formalmente todo el patrimonio, pero excluye selectivamente aquella parte —la empresarial familiar— cuya movilidad y cuyo riesgo de deslocalización son mayores, dejando la carga efectiva del impuesto concentrada, en la práctica, sobre el patrimonio inmobiliario y financiero no afecto a actividad económica.

2.1.4. El esquema técnico del ITSGF

[Seguro] Ya analizamos con detalle, en el subapartado 1.4, la naturaleza del ITSGF como parche armonizador y su blindaje constitucional frente a los recursos autonómicos; falta presentar su diseño técnico propio, porque es una figura normativa autónoma y no un mero apéndice del IP. El hecho imponible es la titularidad, a 31 de diciembre, de un patrimonio neto superior a 3.000.000 de euros, calculado con las mismas reglas de valoración y las mismas exenciones que el IP —incluida, por tanto, la exención de empresa familiar del art. 4.Ocho LIP que desarrollamos en el subapartado 2.1.2—, lo que explica por qué el ITSGF no es, en rigor, un impuesto distinto en su base, sino una cuota adicional sobre la misma base ya depurada por las reglas del IP. [Seguro] Su tarifa reproduce, en sus tramos superiores, la escala estatal del IP —manteniendo el tipo marginal máximo del 3,5%—, pero con un umbral de entrada mucho más elevado que el mínimo exento general del IP, de modo que solo alcanza a los patrimonios verdaderamente cuantiosos. [Seguro] El mecanismo de deducción de la cuota de IP efectivamente satisfecha, ya explicado en el subapartado 1.4, determina que el ITSGF recaude, en la práctica, casi en exclusiva de las Comunidades Autónomas con bonificación total o parcial del IP: los datos de recaudación de 2023 confirman que Madrid, Andalucía y Galicia concentran más del 95% tanto de los contribuyentes como de la recaudación del impuesto, que ascendió ese ejercicio a 623 millones de euros. [Probable] Esta concentración territorial extrema no es un dato accesorio, sino la prueba empírica más directa de que el impuesto cumple exactamente la función que sus impugnantes le atribuían en sus recursos de inconstitucionalidad —neutralizar, allí donde existe, la bonificación autonómica—, con la diferencia de que, tras la STC 149/2023 y la STC 170/2023, esa función queda plenamente amparada por el Tribunal Constitucional como ejercicio legítimo de la potestad tributaria originaria del Estado.


\subsection{Impuesto sobre Bienes Inmuebles}
[Seguro] El IBI se regula, junto con el resto de la fiscalidad municipal, en el Real Decreto Legislativo 2/2004, de 5 de marzo, Texto Refundido de la Ley Reguladora de las Haciendas Locales (TRLRHL), en sus arts. 60 a 77, que lo definen como un tributo directo de carácter real que grava el valor de los bienes inmuebles —urbanos, rústicos y de características especiales— en los términos del propio texto refundido. [Seguro] El hecho imponible no es la propiedad en sentido estricto, sino la titularidad de cualquiera de un conjunto tasado de derechos sobre el inmueble, ordenados jerárquicamente: concesión administrativa sobre el bien o sobre los servicios públicos a que esté afecto; derecho real de superficie; derecho real de usufructo; y, en defecto de los anteriores, derecho de propiedad. [Seguro] La base imponible es, directamente, el valor catastral, fijado por la Dirección General del Catastro —un organismo estatal, no municipal, lo que introduce una divisoria competencial peculiar dentro de un impuesto por lo demás genuinamente local: la valoración es estatal, la gestión tributaria (liquidación, recaudación, aplicación de bonificaciones potestativas) es municipal—, mediante la aprobación periódica de ponencias de valores para cada municipio, que se actualizan mediante coeficientes correctores fijados, típicamente, en las Leyes de Presupuestos Generales del Estado.

2.2.2. La dispersión municipal: tipos de gravamen y la paradoja de la escala

[Probable] Todo este tema ha girado, desde el subapartado 1.4, en torno a la tensión entre un marco formal único y marcos materiales múltiples; el IBI reproduce exactamente esa misma tensión, solo que a una escala que hace palidecer a la autonómica: el art. 72 TRLRHL fija una horquilla legal de tipos de gravamen —aproximadamente entre el 0,4% y el 1,1% para bienes urbanos, con posibilidad de incrementos adicionales en municipios capitales de provincia, que presten servicio de transporte público, o que dispongan de más servicios que los mínimos legalmente exigidos—, dentro de la cual cada uno de los más de 8.000 municipios españoles fija su propio tipo mediante ordenanza fiscal. [Seguro] El resultado es un mapa de dispersión que, en número de unidades territoriales implicadas, supera con creces a las diecisiete Comunidades Autónomas que hemos analizado para IP e ISD: dos inmuebles de valor catastral idéntico pueden soportar una cuota de IBI sustancialmente distinta según el municipio, sin que exista ningún mecanismo de armonización estatal equivalente al ITSGF que analizamos para el IP. [Seguro] A esta dispersión ordinaria se ha sumado, desde la Ley 12/2023, de 24 de mayo, por el derecho a la vivienda —ya citada en el Tema 16 a propósito de las reducciones al arrendador—, la posibilidad de que los ayuntamientos apliquen recargos progresivos sobre viviendas desocupadas: un 50% de recargo sobre la cuota líquida para inmuebles desocupados de forma permanente durante más de dos años, elevable al 100% a partir de tres años, y hasta un 150% para titulares de dos o más inmuebles residenciales desocupados en el mismo término municipal —frente al límite máximo del 50% que regía antes de esta reforma—, con excepciones tasadas (traslado laboral temporal, dependencia, emergencia social). [Suposición] La propia literatura especializada en el sector inmobiliario ha señalado, sin embargo, que el alcance práctico de esta medida es muy limitado: estimaciones basadas en datos de Fotocasa Research sitúan en apenas el 0,1% de los propietarios españoles a quienes, tras aplicar todos los filtros legales (vivienda vacía, más de dos años, sin causa justificada, titularidad múltiple), les resultaría efectivamente aplicable el recargo —un buen ejemplo, para tu exposición oral, de la distancia habitual entre el alcance simbólico-político de una reforma tributaria muy publicitada y su impacto recaudatorio y conductual real, el mismo tipo de contraste que ya documentamos con rigor empírico al hablar de la evaluación de beneficios fiscales en los Temas 16 y 17.

2.2.2. El debate doctrinal de fondo: capacidad económica o principio del beneficio

[Probable] El debate teórico que atraviesa el IBI no es el mismo que hemos desarrollado para el IP: no se discute tanto si debe existir —su legitimidad está mucho menos cuestionada en la literatura comparada que la del impuesto general sobre el patrimonio—, sino qué principio lo fundamenta. La lectura clásica lo sitúa, como al IP, dentro del principio de capacidad económica del art. 31.1 CE: quien es titular de un inmueble tiene, por ese hecho, una manifestación de riqueza gravable. Pero existe una segunda lectura, de raíz distinta, que lo vincula al principio del beneficio: el valor de un inmueble no depende solo de las características físicas del propio bien, sino, en una proporción muy significativa, de la calidad de los servicios públicos municipales que lo rodean —alumbrado, pavimentación, seguridad, colegios, transporte, parques—, de modo que el impuesto sería, en el fondo, una contraprestación por ese beneficio recibido, no una pura medida de capacidad económica personal. [Seguro] Esta segunda lectura tiene una genealogía doctrinal antigua y bien asentada: Henry George, en Progress and Poverty (1879), había defendido ya un "impuesto único" sobre el valor del suelo precisamente con este argumento —el valor de la tierra urbana, a diferencia del valor generado por el trabajo o el capital invertido, procede en gran medida de la localización y de la inversión colectiva de la comunidad circundante, no del esfuerzo individual del propietario, lo que lo convertiría en la base imponible más justa y menos distorsionadora posible—. [Probable] La literatura de hacienda pública local ha sometido esta intuición a contraste empírico bajo el nombre de hipótesis de capitalización: si los servicios públicos locales se "capitalizan" en el valor de los inmuebles —barrios con mejores colegios o menor delincuencia tienen, sistemáticamente, precios inmobiliarios más altos—, entonces gravar el valor del inmueble equivale, indirectamente, a gravar el beneficio recibido de esos servicios, cerrando el círculo con el modelo de Tiebout que ya usamos en el Tema 17: los contribuyentes "votan con los pies" eligiendo municipio según su combinación de impuestos y servicios, y el valor capitalizado en el inmueble revela, precisamente, cuánto están dispuestos a pagar por esa combinación. Wallace Oates, en su trabajo ya clásico (1969, The Effects of Property Taxes and Local Public Spending on Property Values, Journal of Political Economy), encontró evidencia empírica consistente con esta hipótesis de capitalización para municipios de Nueva Jersey, un resultado replicado después en numerosos contextos y que sigue siendo la referencia empírica de cabecera de este debate.

2.2.3. El problema práctico: la desactualización del valor catastral

[Seguro] Más allá del debate de fundamentación, el problema que más ha ocupado a la doctrina y a la práctica española reciente es puramente técnico pero de enorme trascendencia distributiva: los valores catastrales de referencia no se revisan con periodicidad uniforme, sino mediante ponencias de valores aprobadas municipio a municipio, en momentos temporales muy distintos —hay municipios cuya ponencia de valores vigente data de la década de 1984-1990, con un coeficiente de actualización del 1,08 aplicado en 2018, mientras otros la tienen de 2011 o 2012, con coeficientes de actualización inferiores a la unidad, señal de que sus valores estaban ya sobrevalorados respecto del mercado en el momento de fijarlos—. [Probable] La consecuencia es una inequidad horizontal severa entre municipios, e incluso entre barrios de un mismo municipio con distinta antigüedad de revisión: dos inmuebles de valor de mercado idéntico pueden tener bases imponibles de IBI radicalmente distintas según cuándo se aprobó por última vez la ponencia de valores aplicable, sin que exista relación alguna entre esa diferencia y la capacidad económica real de sus titulares. [Suposición] Conviene que distingas con precisión, porque es un punto de confusión habitual incluso en la práctica profesional, que el valor de referencia catastral introducido por la Ley 11/2021 que ya hemos usado repetidamente para ISD, ITPAJD e IP no se aplica al IBI: la propia Exposición de Motivos de esa ley limita expresamente su función a las "operaciones traslativas de inmuebles" —transmisiones onerosas o lucrativas—, mientras que la base imponible del IBI sigue vinculada al valor catastral tradicional, derivado de la ponencia de valores municipal y no de ese nuevo estándar. [Probable] Esta asimetría normativa genera, a su vez, una paradoja poco comentada: un mismo inmueble puede tener, simultáneamente, un valor de referencia (relevante si se transmite, o si su titular es sujeto del IP) sensiblemente distinto de su valor catastral (relevante a efectos de IBI), sin que ambos parámetros, pese a proceder del mismo organismo —la Dirección General del Catastro— y referirse al mismo bien, estén necesariamente coordinados entre sí.

\subsection{Superposición: }

Cierro el apartado retomando la superposición horizontal que anticipamos en el subapartado 1.3. [Probable] Un mismo inmueble no exento —recuerda que la vivienda habitual sí goza de exención parcial en el IP (hasta 300.000 euros) pero no en el IBI, que grava cualquier inmueble con independencia de su condición de vivienda habitual o no— tributa, cada año, dos veces sobre la misma manifestación de riqueza: por IBI, en su condición de bien real, calculado sobre valor catastral y gestionado por el ayuntamiento; y por IP, integrado en el patrimonio neto general del contribuyente, valorado por reglas propias y gestionado por la Comunidad Autónoma o el Estado. 

[Suposición] A diferencia de la superposición IRPF-IP que desarrollamos en el subapartado 1.3, aquí no existe ningún mecanismo legal de corrección o límite conjunto: el legislador no ha considerado que la coexistencia de IBI e IP sobre el mismo inmueble genere el mismo riesgo de confiscatoriedad que sí reconoció, expresamente, para la coexistencia de IRPF e IP —una asimetría de tratamiento que la doctrina no ha explicado con un argumento de principio sólido, y que probablemente se explica mejor por la diferencia de magnitud recaudatoria entre ambas cargas (el IBI, sobre un inmueble concreto, suele representar una fracción muy inferior a la cuota íntegra del propio IP) que por ninguna distinción sustantiva sobre la naturaleza de la superposición.






\clearpage
\section{Imposición sobre el Patrimonio: Transmisión}
\puntosclave{
    
}

Dejamos anunciado, al cerrar el desarrollo histórico del subapartado 1.1, que ISD e ITPAJD no son dos impuestos que coincidan por azar en gravar transmisiones patrimoniales, sino que comparten un tronco normativo común. Cumplo ahora esa promesa con el desarrollo que entonces aplacé.

El Impuesto de Derechos Reales y Transmisión de Bienes —conocido popularmente, incluso mucho después de perder ese nombre oficial, simplemente como "Derechos Reales"— es la figura que, durante la mayor parte del siglo XX español, gravó de forma unificada cualquier cambio de titularidad sobre bienes y derechos, con independencia de que ese cambio se produjera por herencia, donación, compraventa, constitución de un derecho real o cualquier otro negocio jurídico traslativo. [Seguro] Su recorrido normativo está documentado con detalle desde al menos la reforma de 1926 —que introdujo, dentro de esa misma figura unificada, el gravamen sobre el caudal relicto indiviso de la herencia, la modalidad de imposición sucesoria que grava el conjunto del patrimonio hereditario antes de su reparto entre los herederos, frente a la alternativa de gravar las porciones ya adjudicadas a cada uno—, una distinción de diseño que resultará decisiva más adelante en este mismo apartado. [Probable] Durante décadas, "pagar Derechos Reales" fue la expresión coloquial con la que cualquier español designaba, indistintamente, el impuesto sobre una herencia, sobre una donación o sobre la compra de una vivienda de segunda mano —un dato lingüístico que por sí solo revela hasta qué punto estas figuras, hoy separadas, eran percibidas como una sola realidad fiscal—. La escisión definitiva llega ya en democracia: la Ley 32/1980, de 21 de junio, segrega la parte de Derechos Reales relativa a transmisiones onerosas, operaciones societarias y actos documentados, dando origen al actual ITPAJD; y la Ley 29/1987, de 18 de diciembre, segrega la parte relativa a adquisiciones gratuitas —mortis causa e inter vivos—, dando origen al actual ISD. [Probable] Esta genealogía compartida no es un dato erudito prescindible: es la explicación de por qué, como verás en el desarrollo de ambas figuras, comparten todavía hoy mecanismos técnicos comunes —el ya analizado valor de referencia catastral, ciertas reglas de comprobación de valores, e incluso la propia noción de "tributos conexos" que la jurisprudencia reciente ha empezado a reconocer entre ambos impuestos a efectos de plazos de prescripción—. Trato primero el ISD, la figura de mayor densidad doctrinal de todo el tema, y cierro con el ITPAJD.

\subsection{Impuesto sobre Sucesiones y Donaciones}
3.1.1. Encaje y marco normativo

El ISD se regula en la Ley 29/1987, de 18 de diciembre, ya presentada en su origen histórico, desarrollada reglamentariamente por el Real Decreto 1629/1991. [Seguro] Grava tres hechos imponibles distintos, unidos por la nota común de la gratuidad: la adquisición de bienes y derechos por herencia, legado o cualquier otro título sucesorio; la adquisición de bienes y derechos por donación o cualquier otro negocio jurídico gratuito e inter vivos; y la percepción de cantidades por beneficiarios de contratos de seguro de vida cuando el contratante sea persona distinta del beneficiario, salvo que deban tributar por el IRPF. Ya conoces, del Tema 16, la dispersión autonómica y la reducción de empresa familiar de esta figura; no los repito, sino que profundizo ahora en tres cuestiones que entonces no desarrollamos: el debate de diseño técnico del propio impuesto, el gran debate económico sobre su justificación, y el debate específico de su integración con el IRPF.

3.1.2. El debate de diseño: caudal relicto frente a porciones hereditarias

[Seguro] La decisión de diseño más determinante de cualquier impuesto sucesorio comparado —y la que ya anticipamos al hablar de la reforma de 1926— es si el gravamen recae sobre el caudal relicto en su conjunto, antes del reparto entre herederos (el modelo llamado estate tax, propio del sistema anglosajón), o sobre las porciones hereditarias ya adjudicadas a cada heredero individual (el modelo llamado inheritance tax, propio de la tradición continental y del sistema español). [Seguro] El informe de referencia de la OCDE sobre esta materia, Inheritance Taxation in OECD Countries (2021), documenta que de los 24 países de la organización que mantienen imposición sucesoria, la inmensa mayoría —incluidos Alemania, Francia, Italia, Países Bajos, Portugal y España— aplican el modelo de porciones hereditarias, mientras que solo un grupo reducido —Dinamarca, Corea del Sur, Reino Unido y Estados Unidos— mantiene el modelo de caudal relicto. [Probable] La diferencia no es meramente técnica: el modelo de caudal relicto grava una única magnitud —el patrimonio total del causante en el momento del fallecimiento—, con independencia de cuántos herederos lo reciban o de la situación patrimonial de cada uno, lo que simplifica extraordinariamente la gestión (una sola valoración, un solo hecho imponible) pero ignora por completo la capacidad económica individual de cada receptor: bajo este modelo, un patrimonio de un millón de euros tributa exactamente igual si lo hereda un único hijo o si se reparte entre diez. El modelo de porciones hereditarias, en cambio, grava a cada heredero según lo que efectivamente recibe, permitiendo aplicar reducciones y tarifas progresivas calibradas según el grado de parentesco y la cuantía individual adquirida —el propio informe de la OCDE recomienda explícitamente este segundo modelo como preferible desde el punto de vista de la equidad—, pero a costa de una complejidad de gestión muy superior y de un incentivo perverso bien documentado: la posibilidad de fragmentar artificialmente la sucesión entre un número mayor de herederos o de generaciones —incluyendo, por ejemplo, a nietos junto a hijos— precisamente para diluir cada porción individual por debajo de los tramos más gravosos de la tarifa progresiva. [Suposición] España, al adoptar el modelo de porciones, asume conscientemente este segundo riesgo a cambio de la mayor equidad del sistema, y buena parte de la planificación fiscal sucesoria legítima —no elusiva— consiste, precisamente, en estructurar donaciones y particiones de manera que aprovechen al máximo esta arquitectura basada en el receptor individual.

3.1.3. Base imponible, reducciones y tarifa: la arquitectura de la equidad

[Seguro] Coherente con el modelo de porciones que acabamos de fijar, la Ley 29/1987 articula un sistema de reducciones en la base imponible clasificadas por grupos de parentesco —del Grupo I, descendientes menores de 21 años, con las reducciones más generosas, al Grupo IV, extraños y parientes colaterales de cuarto grado en adelante, sin reducción alguna por parentesco—, y aplica después, sobre la cuota resultante, unos coeficientes multiplicadores que la incrementan en función de dos variables simultáneas: el grado de parentesco del heredero con el causante, y el patrimonio preexistente del propio heredero antes de recibir la herencia. [Probable] Este segundo elemento —el coeficiente multiplicador por patrimonio preexistente— es una pieza de diseño particularmente sofisticada y poco conocida fuera del ámbito especializado: significa que la misma herencia, del mismo causante, con el mismo grado de parentesco, tributa más si el heredero que la recibe ya era, antes de heredar, una persona con patrimonio elevado, que si era una persona sin apenas patrimonio previo —una forma de progresividad basada no solo en lo que se recibe, sino en la posición patrimonial de partida de quien lo recibe, que no tiene equivalente exacto en ningún otro impuesto de los que hemos estudiado en este tema—. Ya conoces, del Tema 16, la reducción del 95% por transmisión de empresa familiar (art. 20.2.c) LISD), con requisitos de exención previa en el IP —lo que conecta directamente con la exención del art. 4.Ocho LIP que acabamos de desarrollar con detalle en el apartado 2, y con la misma litigiosidad sobre el alcance del grupo de parentesco que allí documentamos—.

3.1.4. El gran debate económico: ¿por qué gravar —o no gravar— las herencias?

[Probable] Más allá del debate general sobre la imposición patrimonial que desarrollamos en el subapartado 1.2, el ISD tiene su propio debate específico, de una intensidad doctrinal que no tiene parangón en ninguna otra figura de este tema, y que gira en torno a una pregunta empírica previa a cualquier consideración normativa: ¿por qué ahorran las personas para dejar herencia? La respuesta condiciona por completo el efecto que cabe esperar de gravar las herencias sobre el comportamiento económico agregado, y la literatura ofrece, al menos, tres modelos en competencia.

El primero, asociado a Franco Modigliani y a su teoría del ciclo vital del ahorro (formulada originalmente con Albert Ando en 1963), sostiene que el individuo ahorra fundamentalmente para financiar su propio consumo durante la jubilación, no para dejar herencia deliberadamente: cualquier patrimonio que quede al fallecer es, en gran medida, accidental —el resultado de no saber con certeza cuánto se va a vivir, lo que obliga a ahorrar "de más" por precaución frente a una vida más larga de lo esperado—. Bajo este modelo, gravar con dureza las herencias tiene un coste de eficiencia relativamente bajo, porque no está desincentivando un comportamiento de ahorro deliberado: el ahorro se habría producido de todos modos, por motivos de ciclo vital, con independencia de qué ocurra fiscalmente con el remanente al morir.

El segundo modelo, formulado por Laurence Kotlikoff y Lawrence Summers en su artículo ya clásico (1981, The Role of Intergenerational Transfers in Aggregate Capital Accumulation, Journal of Political Economy), llega a una conclusión empírica radicalmente opuesta: usando datos históricos estadounidenses, estiman que el ahorro de ciclo vital no puede explicar más del 20% del stock de capital agregado de Estados Unidos, y que las transferencias intergeneracionales —herencias y donaciones deliberadas— explican la parte dominante, en torno al 80%, de la acumulación de riqueza observada. Bajo su modelo, el ahorro no es principalmente accidental ni motivado por el ciclo vital propio, sino dinástico y deliberado: las personas ahorran, en gran medida, precisamente para dejar herencia a sus descendientes, lo que implica que gravar con dureza las transferencias intergeneracionales sí tendría un efecto significativo sobre la acumulación de capital agregado de la economía —estiman que cada dólar de transferencias reducidas por la vía fiscal reduciría la riqueza agregada en estado estacionario en 70 centavos—. [Seguro] Modigliani replicó con dureza a estas estimaciones (1988, Journal of Economic Perspectives), defendiendo que el ahorro de ciclo vital explica una proporción mucho mayor de la que Kotlikoff y Summers reconocen, en una controversia metodológica —sobre cómo medir correctamente qué parte del patrimonio es "transferido" frente a "generado" por cada generación— que la literatura posterior no ha resuelto de forma concluyente ni en un sentido ni en otro.

[Probable] Existe todavía un tercer modelo, menos conocido pero conceptualmente muy sugerente, formulado por Douglas Bernheim, Andrei Shleifer y el propio Lawrence Summers en un trabajo posterior (1984, Bequests as a Means of Payment, documento de trabajo del NBER que se abre, no por casualidad, citando el verso de El Rey Lear de Shakespeare sobre qué hija ama más a su padre): el legado no sería ni puramente accidental ni puramente altruista, sino estratégico, un instrumento que los padres utilizan deliberadamente para inducir en sus hijos comportamientos deseados —atención, cuidado, visitas— mediante la incertidumbre sobre cómo se repartirá finalmente la herencia. [Suposición] Este tercer modelo, aunque menos citado en el debate de política fiscal que los dos anteriores, tiene una implicación normativa propia y distinta: si el legado es un instrumento de "pago" por servicios de cuidado y atención filial, gravarlo equivale, en cierto sentido, a gravar una transacción de mercado implícita entre generaciones, no una mera transferencia de riqueza sin contrapartida —lo que reabriría, con otro fundamento, el debate sobre si el ISD grava realmente una "ganancia inmerecida" del heredero o algo más próximo a una renta del trabajo de cuidado no formalizada—.

[Probable] Sobre este sustrato empírico incierto se levanta, finalmente, el debate filosófico-político que ya anticipamos en el subapartado 1.2: la tradición rawlsiana, con su exigencia de igualdad equitativa de oportunidades, ve en la sucesión hereditaria una de las fuentes más evidentes de desigualdad de punto de partida entre generaciones, y encuentra en el ISD un instrumento legítimo de corrección. Frente a ella, la tradición libertaria, con Robert Nozick como referencia obligada (Anarchy, State, and Utopia, 1974) y su teoría de la titularidad (entitlement theory), sostiene que si una persona adquirió su patrimonio de forma legítima —tributando ya, en su momento, por la renta que lo generó— y decide transmitirlo voluntariamente a quien libremente elige, ni el donante ni el receptor —ambos actuando de mutuo acuerdo, sin coacción ni externalidad negativa sobre terceros— generan ningún daño que justifique la intervención estatal: gravar esa transferencia sería, bajo esta lectura, una interferencia injustificada en el ejercicio legítimo del derecho de propiedad, particularmente potente como crítica precisamente porque el ISD grava una transacción consentida por ambas partes, a diferencia de casi cualquier otro impuesto, que suele gravar transacciones de mercado sin ese elemento de liberalidad mutua.

3.1.5. El debate de integración: ¿debería el ISD ser, simplemente, IRPF?

[Seguro] Cierro el desarrollo del ISD retomando el ejemplo concreto que planteaste al inicio de este bloque de trabajo: ¿por qué no integrar las adquisiciones gratuitas directamente en el IRPF del receptor, como una renta más, en lugar de mantener un impuesto separado? Ya adelantamos, en el Tema 16, que desde la óptica de SHS una herencia es, sin duda, renta —un incremento del poder de disposición económica del receptor, indistinguible en su naturaleza de cualquier otra ganancia—. [Seguro] El propio informe de la OCDE de 2021 recomienda explícitamente, como buena práctica de diseño, que los impuestos sucesorios se centren en la riqueza recibida por el beneficiario más que en el patrimonio del donante —el modelo de porciones que España ya aplica—, y llega incluso a valorar la opción de integración plena con el impuesto sobre la renta personal, señalando que el argumento de la doble imposición esgrimido en contra de la imposición sucesoria es, en sus propias palabras, "débil, especialmente en el caso de un impuesto basado en el receptor" —precisamente porque quien tributa no es la misma persona que generó y ya tributó por la renta original, sino un tercero distinto que recibe, por primera vez, un incremento patrimonial propio—. [Probable] Frente a este argumento favorable a la integración, la doctrina que defiende mantener el ISD como figura separada esgrime razones de índole tanto técnica como de percepción social. Técnicamente, integrar la herencia en la base general del IRPF del ejercicio en que se recibe generaría picos de progresividad extremos y probablemente injustos —una herencia cuantiosa recibida en un único ejercicio empujaría al heredero, de golpe, a los tramos marginales más altos de la tarifa, sobre una renta que no es recurrente y que no refleja su capacidad económica habitual—, un problema de "moyenne mobile" temporal que cualquier impuesto sobre transacciones puntuales de gran magnitud debe resolver de un modo u otro, y que el ISD resuelve mediante su propia tarifa autónoma, calibrada específicamente para este tipo de adquisición. Desde el punto de vista de la percepción social y de la economía política, la fusión de ambas figuras difuminaría la distinción —relevante para la aceptabilidad pública del tributo, aunque no necesariamente para su fundamento técnico— entre la renta ganada con esfuerzo propio y la renta recibida por mera liberalidad ajena, una distinción que buena parte de la opinión pública considera moralmente relevante incluso cuando los economistas tributarios, siguiendo a SHS, la consideran irrelevante para la medición de la capacidad económica. [Suposición] La razón última, probablemente, por la que ningún país de nuestro entorno ha dado el paso de la integración plena —ni siquiera aquellos, como los países nórdicos, más próximos ideológicamente a maximizar la neutralidad y comprehensividad de sus sistemas fiscales— es de economía política antes que de diseño técnico: un impuesto sucesorio separado, visible y con su propia tarifa, es políticamente más fácil de defender —y de recortar o bonificar selectivamente, como hemos visto en el subapartado 1.4— que una integración que convertiría cualquier reforma del tramo superior del IRPF en, automáticamente, una reforma encubierta de la fiscalidad de las herencias, con la consiguiente pérdida de control político independiente sobre ambas materias.


\subsection{Impuesto sobre Transmisiones Patrimoniales y Actos Jurídicos Documentados}

### 3.2.1. Encaje, marco normativo y la lógica de la triple modalidad

El ITPAJD se regula en el **Real Decreto Legislativo 1/1993, de 24 de septiembre**, cuyo art. 1 —ya lo vimos al principio del tema— configura un tributo **único de naturaleza indirecta** con tres modalidades mutuamente excluyentes: Transmisiones Patrimoniales Onerosas (TPO), Operaciones Societarias (OS) y Actos Jurídicos Documentados (AJD). [Probable] Esta arquitectura tripartita no es casual: responde a tres momentos distintos de la vida económica de un negocio jurídico —la transmisión onerosa de un bien entre particulares (TPO), la constitución o modificación de una sociedad (OS), y la mera formalización documental de un acto ante notario o registrador (AJD)—, y conviene tratarlas por separado porque, como verás, cada una vive hoy una situación radical y desigualmente distinta: TPO sigue siendo el núcleo recaudatorio de la figura; OS ha quedado prácticamente vaciada de contenido por el Derecho de la Unión Europea; y AJD, pese a ser la más "menor" en apariencia —un impuesto sobre el papel notarial—, ha protagonizado el episodio de mayor repercusión jurisprudencial y mediática de todo este tema.

### 3.2.2. TPO: el impuesto residual y la frontera con el IVA

[Seguro] La modalidad TPO grava las transmisiones onerosas de bienes y derechos que no estén sujetas a IVA, así como la constitución de derechos reales, préstamos, fianzas, arrendamientos y concesiones administrativas —recuerda que esta última enumeración procede, casi literalmente, de la vieja Ley 32/1980 que ya citamos al fijar el puente histórico—. [Probable] Su función en el sistema es deliberadamente **residual y complementaria del IVA**: mientras el IVA grava las entregas de bienes y prestaciones de servicios realizadas por empresarios o profesionales en el ejercicio de su actividad, TPO grava exactamente el hueco que el IVA deja libre —fundamentalmente, las transmisiones entre particulares que no actúan como empresarios, de las que la compraventa de vivienda de segunda mano entre dos personas físicas es, con diferencia, el ejemplo de mayor volumen recaudatorio—. [Probable] Esta frontera entre ambos impuestos, sin embargo, no es tan nítida como su formulación teórica sugiere, y genera un problema técnico de calado: el art. 20.Dos de la Ley del IVA permite, en determinados supuestos de segundas y ulteriores entregas de edificaciones —el ámbito natural de la exención de IVA que empuja la operación hacia TPO—, que el transmitente **renuncie a la exención** de IVA cuando el adquirente sea a su vez empresario o profesional con derecho a deducción total del impuesto soportado, optando así, deliberadamente, por tributar por IVA (y, correlativamente, por AJD en lugar de TPO) en vez de por TPO. [Suposición] Esta renuncia no es una curiosidad técnica marginal: es un instrumento de planificación fiscal legítima y muy utilizado en operaciones inmobiliarias entre empresas, precisamente porque el IVA soportado es deducible para el adquirente empresario mientras que TPO, al ser un impuesto sin mecanismo de deducción, se convierte en un coste definitivo —lo que genera un incentivo sistemático a estructurar cualquier operación inmobiliaria de cierto volumen de manera que quede sujeta a IVA en lugar de a TPO siempre que exista la posibilidad legal de hacerlo, y explica por qué buena parte de la litigiosidad práctica de esta modalidad gira, precisamente, en torno a si una operación concreta cumplía o no los requisitos para esa renuncia.

### 3.2.3. OS: la modalidad vaciada por el Derecho de la Unión

[Seguro] La modalidad de Operaciones Societarias grava, en su formulación original, la constitución, el aumento y la disminución de capital, la fusión, la escisión y la disolución de sociedades. [Seguro] Es, de las tres modalidades, la que ha sufrido la transformación más radical en las últimas dos décadas, y por una razón enteramente ajena a la política fiscal española: la **Directiva 2008/7/CE del Consejo, de 12 de febrero de 2008**, relativa a los impuestos indirectos que gravan la concentración de capitales —sustitutiva de la anterior Directiva 69/335/CEE—, limita drásticamente la capacidad de los Estados miembros para gravar las operaciones de aportación de capital a sociedades, en la lógica de que gravar la propia constitución o capitalización de una empresa desincentiva la formación de capital productivo dentro del mercado único europeo, exactamente el mismo argumento de fondo que sustentaba la reserva de capitalización del Impuesto sobre Sociedades que analizamos en el Tema 17. [Seguro] España transpuso esta exigencia mediante la **Ley 4/2008, de 23 de diciembre**, que eliminó la sujeción a OS de las operaciones de constitución de sociedades, aumento de capital, aportaciones de socios que no supongan aumento de capital y traslado de sede de dirección efectiva o domicilio social —vaciando, de facto, el contenido práctico de esta modalidad para la inmensa mayoría de operaciones societarias habituales—. [Probable] El resultado es que, hoy, OS sobrevive casi exclusivamente como modalidad residual para operaciones de disminución de capital con devolución de aportaciones a los socios y para la disolución de sociedades con adjudicación de bienes —un ejemplo elocuente, para tu exposición oral, de cómo una figura tributaria puede permanecer vigente en el texto de la ley y, al mismo tiempo, haber sido prácticamente desactivada en su aplicación real por una norma de armonización europea aprobada con un objetivo completamente distinto del de la propia reforma fiscal española.

### 3.2.4. AJD: naturaleza jurídica discutida y la saga jurisprudencial de 2018

[Probable] La modalidad de Actos Jurídicos Documentados esconde, tras su apariencia de impuesto menor sobre el papel notarial, un debate doctrinal de fondo que rara vez se explicita con claridad: ¿es realmente un **impuesto**, en el sentido técnico de gravar una manifestación de capacidad económica, o es, en su componente de **cuota fija** —el antiguo "timbre" que grava cada folio del documento notarial con independencia de su contenido económico—, más próximo a una **tasa** por la prestación de un servicio de formalización documental? [Suposición] La doctrina tributaria española ha señalado repetidamente esta ambigüedad de naturaleza: la cuota fija no depende en absoluto del valor de lo documentado —dos escrituras de contenido económico radicalmente distinto pagan la misma cuota fija si tienen el mismo número de folios—, lo que la aleja del principio de capacidad económica del art. 31.1 CE y la acerca conceptualmente a una tasa por servicios registrales o notariales; la **cuota variable o gradual**, en cambio —la que grava, con un tipo ad valorem, los documentos notariales que tengan por objeto cantidad o cosa valuable, contengan actos inscribibles en un registro público, y no estén sujetos a ISD ni a las otras dos modalidades del propio ITPAJD—, sí opera como un auténtico impuesto sobre una magnitud económica, y es precisamente esta cuota variable la que protagoniza el episodio jurisprudencial más relevante de toda esta figura.

[Seguro] El sujeto pasivo de la cuota variable del AJD en las escrituras de constitución de préstamo con garantía hipotecaria fue, durante décadas, una cuestión pacífica: el art. 29 del RDLeg 1/1993, interpretado en conexión con el art. 68.2 de su Reglamento, designaba como sujeto pasivo al **adquirente del bien o derecho** documentado, y la jurisprudencia consolidada de la Sala Tercera del Tribunal Supremo entendía que ese "adquirente" era, en una escritura de préstamo hipotecario, el **prestatario** —el cliente que recibía el préstamo y constituía la hipoteca como garantía—, no el banco prestamista. [Seguro] Esta doctrina asentada saltó por los aires el **16 de octubre de 2018**, cuando la Sección Segunda de la Sala Tercera del Tribunal Supremo (Sentencia núm. 1505/2018), en un giro jurisprudencial radical, declaró que el verdadero sujeto pasivo era la **entidad prestamista** —razonando que es el banco quien tiene el interés jurídico y económico primordial en que el documento se formalice e inscriba, porque es la inscripción registral de la hipoteca la que le proporciona la garantía real que hace valioso el préstamo concedido—, y **anuló, por ilegal**, el art. 68.2 del Reglamento del impuesto, que hasta entonces avalaba la interpretación contraria. Esta sentencia fue confirmada, en los días siguientes, por otras dos resoluciones de la misma Sección, de 22 y 23 de octubre de 2018.

[Seguro] La reacción institucional fue inmediata y de una magnitud excepcional: el propio Presidente de la Sala Tercera, ante la trascendencia económica y social del cambio de criterio —afectaba, con efectos potencialmente retroactivos, a millones de hipotecas ya formalizadas—, convocó al **Pleno de la Sala** para reconsiderar la cuestión, algo infrecuente para una sala que ya se había pronunciado en tres sentencias consecutivas. El Pleno, reunido el **6 de noviembre de 2018**, y tras una votación reñidísima —**15 votos frente a 13**—, decidió **revertir** el criterio de las tres sentencias de octubre, volviendo a la doctrina tradicional de que el sujeto pasivo es el prestatario. [Seguro] Ante la inseguridad jurídica generada por esta sucesión de vaivenes —tres sentencias en un sentido, después el propio Pleno de la misma Sala en el sentido contrario, en apenas tres semanas—, el Gobierno reaccionó por la vía más rápida posible: el **Real Decreto-ley 17/2018, de 8 de noviembre**, publicado en el BOE precisamente el mismo día en que se hacían públicas las sentencias de octubre que ya habían quedado superadas por el Pleno, modificó los arts. 29 y 45.I.B) del RDLeg 1/1993 para atribuir, con efectos desde el 10 de noviembre de 2018 y **hacia el futuro**, la condición de sujeto pasivo de la cuota variable de AJD en préstamos hipotecarios al **prestamista**, cerrando definitivamente por vía legislativa una controversia que la propia jurisprudencia no había conseguido estabilizar. [Probable] El episodio es, más allá de su enorme repercusión mediática y de su interés como "gancho" expositivo evidente, un caso de estudio sobre el **principio de seguridad jurídica** del art. 9.3 CE en su vertiente jurisdiccional: la doctrina especializada que analizó el episodio en 2019 señaló que la secuencia completa —cambio de doctrina consolidada, reversión por el propio tribunal en cuestión de semanas, e intervención legislativa urgente para zanjar la incertidumbre— constituye un ejemplo paradigmático de cómo la inestabilidad jurisprudencial sobre una cuestión de alcance masivo puede generar, por sí misma, un problema de seguridad jurídica tan grave como el que pretendía resolver el cambio de criterio original, obligando al legislador a intervenir no tanto para corregir una interpretación que considerara errónea, sino simplemente para **poner fin a la incertidumbre** sobre cuál era el derecho aplicable en un momento dado.




\subsection{Superposición intertemporal:}
Queda pendiente la tercera manifestación de superposición que anunciamos en el subapartado 1.3: la que conecta, no dos figuras que actúan simultáneamente sobre el mismo bien, sino una figura de transmisión con las figuras de tenencia que se activan después de ella. Conviene precisar, antes de desarrollarla, que no es, en rigor, una superposición en el mismo sentido técnico que la del IRPF-IP o la del IP-IBI que ya analizamos —no hay aquí dos hechos imponibles distintos gravando simultáneamente la misma magnitud—, sino una secuencia: el ISD grava la adquisición, un hecho imponible instantáneo y puntual; el IP y el IBI gravan la tenencia posterior de lo adquirido, hechos imponibles periódicos y recurrentes. Son, conceptualmente, momentos distintos del ciclo vital del mismo bien, no el mismo momento gravado dos veces.

[Probable] Pero esta distinción conceptual no elimina la relevancia económica real de la secuencia, y es aquí donde conviene que cierres el tema con una síntesis que conecte, de forma explícita, con el Tema 16. Un heredero que acepta una herencia no solo asume la cuota del ISD correspondiente a la adquisición: asume también, desde ese mismo momento, un flujo de obligaciones tributarias futuras —IP anual si su patrimonio total supera el mínimo exento, IBI anual sobre cualquier inmueble recibido— que convierte el coste fiscal real de aceptar una herencia en algo sustancialmente superior al que refleja, por sí sola, la liquidación del ISD. [Suposición] Este dato tiene una implicación práctica poco comentada fuera del ámbito especializado: puede generar un incentivo económico real hacia la liquidación inmediata del patrimonio heredado —vender el inmueble o la cartera de valores poco después de heredar, en lugar de conservarlos— antes que hacia su mantenimiento a largo plazo, precisamente porque vender pronto evita el flujo futuro de IP e IBI, mientras que conservar el bien lo perpetúa. Y aquí es donde el tema se cierra sobre sí mismo con una conexión directa al Tema 16: recordarás que el art. 36 LIRPF, en relación con el art. 33.3.b), determina que el heredero que vende poco después de heredar parte de un valor de adquisición actualizado al momento de la transmisión mortis causa —la llamada "plusvalía del muerto" no sujeta—, de modo que la venta inmediata genera, con frecuencia, poca o ninguna ganancia patrimonial gravable en IRPF. [Suposición] El resultado agregado de este conjunto de reglas —ISD en la adquisición, exención de la plusvalía acumulada hasta la muerte, ausencia de plusvalía relevante en una venta rápida posterior, y evitación del flujo futuro de IP/IBI mediante esa misma venta— es que el sistema tributario español, tomado en su conjunto, no es neutral respecto de la decisión de conservar o liquidar un patrimonio heredado: favorece, en términos puramente fiscales, la liquidación temprana frente a la continuidad patrimonial familiar a largo plazo, salvo, precisamente, en el caso de la empresa familiar, donde la reducción del 95% en ISD y la exención en IP —ambas condicionadas, como vimos, al mantenimiento de la actividad— invierten deliberadamente ese incentivo. Es una buena síntesis de cierre, porque muestra que las cuatro figuras del tema, lejos de operar en compartimentos estancos, configuran entre sí —y con el propio IRPF del Tema 16— un sistema de incentivos coherente, aunque no siempre explícitamente diseñado como tal.





































































\clearpage
\section{Conclusión} 


\subsection{Recapitulación}


\subsection{Extensiones y relación con otras partes del Temario}



\subsection{Opinión}

\subsubsection{Idea final (Salida o cierra)}

\clearpage
\pagenumbering{gobble}
\MyBibliography

SIMONS, H. (1938), ya citado en el Tema 16; ahora en su proyección sobre la capacidad económica del patrimonio como magnitud autónoma de la renta.
PIKETTY, T. (2013): *Le Capital au XXIe siècle*, Éditions du Seuil, París (edición en castellano: *El capital en el siglo XXI*, Fondo de Cultura Económica).
ZUCMAN, G. (2015): *The Hidden Wealth of Nations: The Scourge of Tax Havens*, University of Chicago Press.
ZUCMAN, G. (2024): *A Blueprint for a Coordinated Minimum Effective Taxation Standard for Ultra-High-Net-Worth Individuals*, informe encargado por la Presidencia brasileña del G20, junio de 2024.
SAEZ, E. y ZUCMAN, G. (2019): "Response to Summers and Sarin, 'A Wealth Tax Presents a Revenue Estimation Puzzle'", 25 de junio de 2019.
SARIN, N. y SUMMERS, L. H. (2019): "A 'Wealth Tax' Presents a Revenue Estimation Puzzle", *The Washington Post*, 4 de abril de 2019, y "Be Very Skeptical About How Much Revenue Elizabeth Warren's Wealth Tax Could Generate", *The Washington Post*, 28 de junio de 2019.
Ley 19/1991, de 6 de junio, del Impuesto sobre el Patrimonio, art. 31.Uno (ya citada; ahora en su contenido específico sobre el límite conjunto IRPF-IP).
Convenio Europeo de Derechos Humanos, Protocolo Adicional 1, art. 1 (protección de la propiedad, límite a la confiscatoriedad fiscal según jurisprudencia del TEDH).
Tratado de Funcionamiento de la Unión Europea, art. 63 (libre circulación de capitales).
Sentencias del Tribunal Supremo de 29 de octubre de 2025 (rec. 4701/2023) y de 3 de noviembre de 2025 (rec. 7626/2023), Sala de lo Contencioso-Administrativo.
DE JUAN, J. (2024): comentario en *Crónica Tributaria* sobre la compatibilidad del límite conjunto IRPF-IP con la jurisprudencia del TEDH sobre confiscatoriedad.
- Ley 38/2022, de 27 de diciembre, para el establecimiento de gravámenes temporales energético y de entidades de crédito y establecimientos financieros de crédito, y por la que se crea el Impuesto Temporal de Solidaridad de las Grandes Fortunas (art. 3).
- Ley 21/2001, de 27 de diciembre (ya citada en el subapartado 1.1; ahora en su contenido específico de capacidad normativa cedida sobre IP, ISD e ITPAJD).
- Ley Orgánica 8/1980 (LOFCA), art. 2.1.g) (principio de lealtad institucional).
- Sentencia del Tribunal Constitucional 149/2023, de 7 de noviembre (recurso de inconstitucionalidad de la Comunidad de Madrid).
- Sentencia del Tribunal Constitucional 170/2023, de 22 de noviembre (recurso de inconstitucionalidad de la Junta de Andalucía).
Ley 19/1991, de 6 de junio, del Impuesto sobre el Patrimonio, arts. 4.Ocho, 9 y 10 (ya citada; ahora en su contenido específico).
Consulta Vinculante de la DGT V2390-23, de 5 de septiembre de 2023 (grupo de parentesco y funciones de dirección).
Consulta Vinculante de la DGT V0426-25, de 20 de marzo de 2025 (exclusión de sobrinos del grupo de parentesco).
Sentencias del Tribunal Supremo 1776/2016 y 1198/2016 (basta que las funciones de dirección las ejerza cualquier miembro del grupo de parentesco).
Régimen del Impôt sur la Fortune Immobilière (IFI) francés, creado por la Loi de finances pour 2018 (Francia), como término de comparación del modelo de imposición patrimonial limitada a inmuebles.
Real Decreto Legislativo 2/2004, de 5 de marzo, Texto Refundido de la Ley Reguladora de las Haciendas Locales, arts. 60 a 77.
Real Decreto Legislativo 1/2004, de 5 de marzo, Texto Refundido de la Ley del Catastro Inmobiliario (relación entre valor catastral y valor de referencia).
GEORGE, H. (1879): Progress and Poverty, Nueva York (formulación del impuesto único sobre el valor del suelo).
OATES, W. E. (1969): "The Effects of Property Taxes and Local Public Spending on Property Values: An Empirical Study of Tax Capitalization and the Tiebout Hypothesis", Journal of Political Economy, vol. 77, núm. 6.
OCDE (2021): Inheritance Taxation in OECD Countries, OECD Tax Policy Studies, París.
KOTLIKOFF, L. J. y SUMMERS, L. H. (1981): "The Role of Intergenerational Transfers in Aggregate Capital Accumulation", Journal of Political Economy, vol. 89, núm. 4.
MODIGLIANI, F. (1988): "The Role of Intergenerational Transfers and Life Cycle Saving in the Accumulation of Wealth", Journal of Economic Perspectives, vol. 2, núm. 2.
BERNHEIM, B. D., SHLEIFER, A. y SUMMERS, L. H. (1984): "Bequests as a Means of Payment", NBER Working Paper 1303.
NOZICK, R. (1974): Anarchy, State, and Utopia, Basic Books, Nueva York.
ANDO, A. y MODIGLIANI, F. (1963): "The 'Life Cycle' Hypothesis of Saving: Aggregate Implications and Tests", American Economic Review, vol. 53, núm. 1.
- Ley 37/1992, de 28 de diciembre, del Impuesto sobre el Valor Añadido, art. 20.Dos (renuncia a la exención en segundas entregas de edificaciones).
- Directiva 2008/7/CE del Consejo, de 12 de febrero de 2008, relativa a los impuestos indirectos que gravan la concentración de capitales.
- Ley 4/2008, de 23 de diciembre, por la que se suprime el gravamen del Impuesto sobre el Patrimonio, se generaliza el sistema de devolución mensual en el IVA, y se introducen otras modificaciones en la normativa tributaria (exención de operaciones societarias).
- Sentencia del Tribunal Supremo 1505/2018, de 16 de octubre de 2018, Sala de lo Contencioso-Administrativo, Sección Segunda.
- Acuerdo del Pleno de la Sala Tercera del Tribunal Supremo de 6 de noviembre de 2018.
- Real Decreto-ley 17/2018, de 8 de noviembre, por el que se modifica el Texto Refundido de la Ley del Impuesto sobre Transmisiones Patrimoniales y Actos Jurídicos Documentados.
- ADAME MARTÍNEZ, F. (2018): "El sujeto pasivo del Impuesto sobre Actos Jurídicos Documentados en los préstamos con garantía hipotecaria", *Revista de Contabilidad y Tributación. CEF*, núm. 429.




\MyBibItem{DAWID y GATTI}{Chapter 2 - Agent-Based Macroeconomics}{2018}{https://doi.org/10.1016/bs.hescom.2018.02.006}


% ANEXOS
\clearpage










\end{document}